
\section{深度学习，真的行}



\section{线性神经网络}

在介绍深度神经网络之前，我们需要了解神经网络训练的基础知识。本章我们将介绍神经网络的整个训练过程，包括：定义简单的神经网络架构、数据处理、指定损失函数和如何训练模型。为了更容易学习，我们将从经典算法------线性神经网络开始，介绍神经网络的基础知识。经典统计学习技术中的线性回归和softmax回归可以视为线性神经网络，这些知识将为本书其他部分中更复杂的技术奠定基础。

-3.1.线性回归
%-[3.1.1.线性回归的基本元素](https://zh-v2.d2l.ai/chapter_linear-networks/linear-regression.html#id2)
%-[3.1.2.矢量化加速](https://zh-v2.d2l.ai/chapter_linear-networks/linear-regression.html#id8)
%-[3.1.3.正态分布与平方损失](https://zh-v2.d2l.ai/chapter_linear-networks/linear-regression.html#subsec-normal-distribution-and-squared-loss)
%-[3.1.4.从线性回归到深度网络](https://zh-v2.d2l.ai/chapter_linear-networks/linear-regression.html#id10)
%-[3.1.5.小结](https://zh-v2.d2l.ai/chapter_linear-networks/linear-regression.html#id14)
%-[3.1.6.练习](https://zh-v2.d2l.ai/chapter_linear-networks/linear-regression.html#id15)
%-3.2.线性回归的从零开始实现
%-[3.2.1.生成数据集](https://zh-v2.d2l.ai/chapter_linear-networks/linear-regression-scratch.html#id2)
%-[3.2.2.读取数据集](https://zh-v2.d2l.ai/chapter_linear-networks/linear-regression-scratch.html#id3)
%-[3.2.3.初始化模型参数](https://zh-v2.d2l.ai/chapter_linear-networks/linear-regression-scratch.html#id4)
%-[3.2.4.定义模型](https://zh-v2.d2l.ai/chapter_linear-networks/linear-regression-scratch.html#id5)
%-[3.2.5.定义损失函数](https://zh-v2.d2l.ai/chapter_linear-networks/linear-regression-scratch.html#id6)
%-[3.2.6.定义优化算法](https://zh-v2.d2l.ai/chapter_linear-networks/linear-regression-scratch.html#id7)
%-[3.2.7.训练](https://zh-v2.d2l.ai/chapter_linear-networks/linear-regression-scratch.html#id8)
%-[3.2.8.小结](https://zh-v2.d2l.ai/chapter_linear-networks/linear-regression-scratch.html#id9)
%-[3.2.9.练习](https://zh-v2.d2l.ai/chapter_linear-networks/linear-regression-scratch.html#id10)
%-3.3.线性回归的简洁实现
%-[3.3.1.生成数据集](https://zh-v2.d2l.ai/chapter_linear-networks/linear-regression-concise.html#id2)
%-[3.3.2.读取数据集](https://zh-v2.d2l.ai/chapter_linear-networks/linear-regression-concise.html#id3)
%-[3.3.3.定义模型](https://zh-v2.d2l.ai/chapter_linear-networks/linear-regression-concise.html#id4)
%-[3.3.4.初始化模型参数](https://zh-v2.d2l.ai/chapter_linear-networks/linear-regression-concise.html#id5)
%-[3.3.5.定义损失函数](https://zh-v2.d2l.ai/chapter_linear-networks/linear-regression-concise.html#id6)
%-[3.3.6.定义优化算法](https://zh-v2.d2l.ai/chapter_linear-networks/linear-regression-concise.html#id7)
%-[3.3.7.训练](https://zh-v2.d2l.ai/chapter_linear-networks/linear-regression-concise.html#id8)
%-[3.3.8.小结](https://zh-v2.d2l.ai/chapter_linear-networks/linear-regression-concise.html#id9)
%-[3.3.9.练习](https://zh-v2.d2l.ai/chapter_linear-networks/linear-regression-concise.html#id10)
%-3.4.softmax回归
%-[3.4.1.分类问题](https://zh-v2.d2l.ai/chapter_linear-networks/softmax-regression.html#subsec-classification-problem)
%-[3.4.2.网络架构](https://zh-v2.d2l.ai/chapter_linear-networks/softmax-regression.html#id2)
%-[3.4.3.全连接层的参数开销](https://zh-v2.d2l.ai/chapter_linear-networks/softmax-regression.html#subsec-parameterization-cost-fc-layers)
%-[3.4.4.softmax运算](https://zh-v2.d2l.ai/chapter_linear-networks/softmax-regression.html#subsec-softmax-operation)
%-[3.4.5.小批量样本的矢量化](https://zh-v2.d2l.ai/chapter_linear-networks/softmax-regression.html#subsec-softmax-vectorization)
%-[3.4.6.损失函数](https://zh-v2.d2l.ai/chapter_linear-networks/softmax-regression.html#id7)
%-[3.4.7.信息论基础](https://zh-v2.d2l.ai/chapter_linear-networks/softmax-regression.html#subsec-info-theory-basics)
%-[3.4.8.模型预测和评估](https://zh-v2.d2l.ai/chapter_linear-networks/softmax-regression.html#id15)
%-[3.4.9.小结](https://zh-v2.d2l.ai/chapter_linear-networks/softmax-regression.html#id16)
%-[3.4.10.练习](https://zh-v2.d2l.ai/chapter_linear-networks/softmax-regression.html#id17)
%-3.5.图像分类数据集
%-[3.5.1.读取数据集](https://zh-v2.d2l.ai/chapter_linear-networks/image-classification-dataset.html#id4)
%-[3.5.2.读取小批量](https://zh-v2.d2l.ai/chapter_linear-networks/image-classification-dataset.html#id5)
%-[3.5.3.整合所有组件](https://zh-v2.d2l.ai/chapter_linear-networks/image-classification-dataset.html#id6)
%-[3.5.4.小结](https://zh-v2.d2l.ai/chapter_linear-networks/image-classification-dataset.html#id7)
%-[3.5.5.练习](https://zh-v2.d2l.ai/chapter_linear-networks/image-classification-dataset.html#id8)
%-3.6.softmax回归的从零开始实现
%-[3.6.1.初始化模型参数](https://zh-v2.d2l.ai/chapter_linear-networks/softmax-regression-scratch.html#id1)
%-[3.6.2.定义softmax操作](https://zh-v2.d2l.ai/chapter_linear-networks/softmax-regression-scratch.html#id2)
%-[3.6.3.定义模型](https://zh-v2.d2l.ai/chapter_linear-networks/softmax-regression-scratch.html#id3)
%-[3.6.4.定义损失函数](https://zh-v2.d2l.ai/chapter_linear-networks/softmax-regression-scratch.html#id4)
%-[3.6.5.分类精度](https://zh-v2.d2l.ai/chapter_linear-networks/softmax-regression-scratch.html#id5)
%-[3.6.6.训练](https://zh-v2.d2l.ai/chapter_linear-networks/softmax-regression-scratch.html#id6)
%-[3.6.7.预测](https://zh-v2.d2l.ai/chapter_linear-networks/softmax-regression-scratch.html#id7)
%-[3.6.8.小结](https://zh-v2.d2l.ai/chapter_linear-networks/softmax-regression-scratch.html#id8)
%-[3.6.9.练习](https://zh-v2.d2l.ai/chapter_linear-networks/softmax-regression-scratch.html#id9)
%-3.7.softmax回归的简洁实现
%-[3.7.1.初始化模型参数](https://zh-v2.d2l.ai/chapter_linear-networks/softmax-regression-concise.html#id1)
%-[3.7.2.重新审视Softmax的实现](https://zh-v2.d2l.ai/chapter_linear-networks/softmax-regression-concise.html#subsec-softmax-implementation-revisited)
%-[3.7.3.优化算法](https://zh-v2.d2l.ai/chapter_linear-networks/softmax-regression-concise.html#id3)
%-[3.7.4.训练](https://zh-v2.d2l.ai/chapter_linear-networks/softmax-regression-concise.html#id4)
%-[3.7.5.小结](https://zh-v2.d2l.ai/chapter_linear-networks/softmax-regression-concise.html#id5)
%-[3.7.6.练习](https://zh-v2.d2l.ai/chapter_linear-networks/softmax-regression-concise.html#id6)



通俗来说，机器学习是一门讨论各式各样的适用于不同问题的函数形式，以及如何使用数据来有效地获取函数参数具体值的学科。深度学习是指机器学习中的一类函数，它们的形式通常为多层神经网络。近年来，仰仗着大数据集和强大的硬件，深度学习已逐渐成为处理图像、文本语料和声音信号等复杂高维度数据的主要方法。

我们现在正处于一个程序设计得到深度学习的帮助越来越多的时代。这可以说是计算机科学历史上的一个分水岭。举个例子，深度学习已经在你的手机里：拼写校正、语音识别、认出社交媒体照片里的好友们等。得益于优秀的算法、快速而廉价的算力、前所未有的大量数据以及强大的软件工具，如今大多数软件工程师都有能力建立复杂的模型来解决十年前连最优秀的科学家都觉得棘手的问题。

\subsection{.1.起源}

虽然深度学习似乎是最近几年刚兴起的名词，但它所基于的神经网络模型和用数据编程的核心思想已经被研究了数百年。自古以来，人类就一直渴望能从数据中分析出预知未来的窍门。实际上，数据分析正是大部分自然科学的本质，我们希望从日常的观测中提取规则，并找寻不确定性。

早在17世纪，[雅各比·伯努利（1655–1705）]%(https：//en.wikipedia.org/wiki/Jacob_Bernoulli)
提出了描述只有两种结果的随机过程（如抛掷一枚硬币）的伯努利分布。大约一个世纪之后，[卡尔·弗里德里希·高斯（1777–1855）]%(https：//en.wikipedia.org/wiki/Carl_Friedrich_Gauss)
发明了今日仍广泛用在从保险计算到医学诊断等领域的最小二乘法。概率论、统计学和模式识别等工具帮助自然科学的实验学家们从数据回归到自然定律，从而发现了如欧姆定律（描述电阻两端电压和流经电阻电流关系的定律）这类可以用线性模型完美表达的一系列自然法则。

即使是在中世纪，数学家也热衷于利用统计学来做出估计。例如，在[雅各比·科贝尔（1460–1533）]
%(https：//www.maa.org/press/periodicals/convergence/mathematical-treasures-jacob-kobels-geometry)
的几何书中记载了使用16名男子的平均脚长来估计男子的平均脚长。


现代统计学在20世纪的真正起飞要归功于数据的收集和发布。统计学巨匠之一[罗纳德·费雪（1890–1962）]
%(https：//en.wikipedia.org/wiki/Ronald_Fisher)
对统计学理论和统计学在基因学中的应用功不可没。他发明的许多算法和公式，例如线性判别分析和费雪信息，仍经常被使用。即使是他在1936年发布的Iris数据集，仍然偶尔被用于演示机器学习算法。

[克劳德·香农（1916–2001）]%(https：//en.wikipedia.org/wiki/Claude_Shannon)
的信息论以及[阿兰·图灵（1912–1954）]%(https：//en.wikipedia.org/wiki/Allan_Turing)
的计算理论也对机器学习有深远影响。图灵在他著名的论文[《计算机器与智能》]
%(https：//www.jstor.org/stable/2251299)
中提出了“机器可以思考吗？”这样一个问题[1]。在他描述的“图灵测试”中，如果一个人在使用文本交互时不能区分他的对话对象到底是人类还是机器的话，那么即可认为这台机器是有智能的。时至今日，智能机器的发展可谓日新月异。

另一个对深度学习有重大影响的领域是神经科学与心理学。既然人类显然能够展现出智能，那么对于解释并逆向工程人类智能机理的探究也在情理之中。最早的算法之一是由[唐纳德·赫布（1904–1985）]%(https：//en.wikipedia.org/wiki/Donald_O._Hebb)
正式提出的。在他开创性的著作[《行为的组织》]
%(http：//s-f-walker.org.uk/pubsebooks/pdfs/The_Organization_of_Behavior-Donald_O._Hebb.pdf)
中，他提出神经是通过正向强化来学习的，即赫布理论[2]。赫布理论是感知机学习算法的原型，并成为支撑今日深度学习的随机梯度下降算法的基石：强化合意的行为、惩罚不合意的行为，最终获得优良的神经网络参数。

来源于生物学的灵感是神经网络名字的由来。这类研究者可以追溯到一个多世纪前的[亚历山大·贝恩（1818–1903）]
%(https：//en.wikipedia.org/wiki/Alexander_Bain)
和[查尔斯·斯科特·谢灵顿（1857–1952）]
%(https：//en.wikipedia.org/wiki/Charles_Scott_Sherrington)
。研究者们尝试组建模仿神经元互动的计算电路。随着时间发展，神经网络的生物学解释被稀释，但仍保留了这个名字。时至今日，绝大多数神经网络都包含以下的核心原则。

-交替使用线性处理单元与非线性处理单元，它们经常被称为“层”。
-使用链式法则（即反向传播）来更新网络的参数。

在最初的快速发展之后，自约1995年起至2005年，大部分机器学习研究者的视线从神经网络上移开了。这是由于多种原因。首先，训练神经网络需要极强的计算力。尽管20世纪末内存已经足够，计算力却不够充足。其次，当时使用的数据集也相对小得多。费雪在1936年发布的的Iris数据集仅有150个样本，并被广泛用于测试算法的性能。具有6万个样本的MNIST数据集在当时已经被认为是非常庞大了，尽管它如今已被认为是典型的简单数据集。由于数据和计算力的稀缺，从经验上来说，如核方法、决策树和概率图模型等统计工具更优。它们不像神经网络一样需要长时间的训练，并且在强大的理论保证下提供可以预测的结果。

\subsection{1.2.发展}

互联网的崛起、价廉物美的传感器和低价的存储器令我们越来越容易获取大量数据。加之便宜的计算力，尤其是原本为电脑游戏设计的GPU的出现，上文描述的情况改变了许多。一瞬间，原本被认为不可能的算法和模型变得触手可及。这样的发展趋势从如下表格中可见一斑。

\begin{table}[h]
\centering
\begin{tabular}{|c|c|c|c|}
\hline
\textbf{年代} & \textbf{数据样本个数} & \textbf{内存} & \textbf{每秒浮点计算数} \\
\hline
1970 & $10^2$ （Iris） & 1 KB & 100 K（Intel 8080） \\
1980 & $10^3$ （波士顿房价） & 100 KB & 1 M（Intel 80186） \\
1990 & $10^4$ （手写字符识别） & 10 MB & 10 M（Intel 80486） \\
2000 & $10^7$ （网页） & 100 MB & 1 G（Intel Core） \\
2010 & $10^{10}$ （广告） & 1 GB & 1 T（NVIDIA C2050） \\
2020 & $10^{12}$ （社交网络） & 100 GB & 1 P（NVIDIA DGX-2） \\
\hline
\end{tabular}
\caption{机器学习发展历程中的数据规模、内存和计算力变化}
\label{tab:ml_evolution}
\end{table}

很显然，存储容量没能跟上数据量增长的步伐。与此同时，计算力的增长又盖过了数据量的增长。这样的趋势使得统计模型可以在优化参数上投入更多的计算力，但同时需要提高存储的利用效率，例如使用非线性处理单元。这也相应导致了机器学习和统计学的最优选择从广义线性模型及核方法变化为深度多层神经网络。这样的变化正是诸如多层感知机、卷积神经网络、长短期记忆循环神经网络和Q学习等深度学习的支柱模型在过去10年从坐了数十年的冷板凳上站起来被“重新发现”的原因。

近年来在统计模型、应用和算法上的进展常被拿来与寒武纪大爆发（历史上物种数量大爆发的一个时期）做比较。但这些进展不仅仅是因为可用资源变多了而让我们得以用新瓶装旧酒。下面的列表仅仅涵盖了近十年来深度学习长足发展的部分原因。

-优秀的容量控制方法，如丢弃法，使大型网络的训练不再受制于过拟合（大型神经网络学会记忆大部分训练数据的行为）[3]。这是靠在整个网络中注入噪声而达到的，如训练时随机将权重替换为随机的数字[4]。
-注意力机制解决了另一个困扰统计学超过一个世纪的问题：如何在不增加参数的情况下扩展一个系统的记忆容量和复杂度。注意力机制使用了一个可学习的指针结构来构建出一个精妙的解决方法[5]。也就是说，与其在像机器翻译这样的任务中记忆整个句子，不如记忆指向翻译的中间状态的指针。由于生成译文前不需要再存储整句原文的信息，这样的结构使准确翻译长句变得可能。
-记忆网络[6]和神经编码器—解释器[7]这样的多阶设计使得针对推理过程的迭代建模方法变得可能。这些模型允许重复修改深度网络的内部状态，这样就能模拟出推理链条上的各个步骤，就好像处理器在计算过程中修改内存一样。
-另一个重大发展是生成对抗网络的发明[8]。传统上，用在概率分布估计和生成模型上的统计方法更多地关注于找寻正确的概率分布，以及正确的采样算法。生成对抗网络的关键创新在于将采样部分替换成了任意的含有可微分参数的算法。这些参数将被训练到使辨别器不能再分辨真实的和生成的样本。生成对抗网络可使用任意算法来生成输出的这一特性为许多技巧打开了新的大门。例如生成奔跑的斑马[9]和生成名流的照片[10]都是生成对抗网络发展的见证。
-许多情况下单块GPU已经不能满足在大型数据集上进行训练的需要。过去10年内我们构建分布式并行训练算法的能力已经有了极大的提升。设计可扩展算法的最大瓶颈在于深度学习优化算法的核心：随机梯度下降需要相对更小的批量。与此同时，更小的批量也会降低GPU的效率。如果使用1，024块GPU，每块GPU的批量大小为32个样本，那么单步训练的批量大小将是32，000个以上。近年来李沐[11]、YangYou等人[12]以及XianyanJia等人[13]的工作将批量大小增至多达64，000个样例，并把在ImageNet数据集上训练ResNet-50模型的时间降到了7分钟。与之相比，最初的训练时间需要以天来计算。
-并行计算的能力也为至少在可以采用模拟情况下的强化学习的发展贡献了力量。并行计算帮助计算机在围棋、雅达利游戏、星际争霸和物理模拟上达到了超过人类的水准。
-深度学习框架也在传播深度学习思想的过程中扮演了重要角色。[Caffe](https：//github.com/BVLC/caffe)、[Torch](https：//github.com/torch)和[Theano](https：//github.com/Theano/Theano)这样的第一代框架使建模变得更简单。许多开创性的论文都用到了这些框架。如今它们已经被[TensorFlow](https：//github.com/tensorflow/tensorflow)（经常是以高层API[Keras](https：//github.com/keras-team/keras)的形式被使用）、[CNTK](https：//github.com/Microsoft/CNTK)、[Caffe2](https：//github.com/caffe2/caffe2)和[ApacheMXNet](https：//github.com/apache/incubator-mxnet)所取代。第三代，即命令式深度学习框架，是由用类似NumPy的语法来定义模型的[Chainer](https：//github.com/chainer/chainer)所开创的。这样的思想后来被[PyTorch](https：//github.com/pytorch/pytorch)和MXNet的[GluonAPI](https：//github.com/apache/incubator-mxnet)采用，后者也正是本书用来教学深度学习的工具。

系统研究者负责构建更好的工具，统计学家建立更好的模型。这样的分工使工作大大简化。举例来说，在2014年时，训练一个逻辑回归模型曾是卡内基梅隆大学布置给机器学习方向的新入学博士生的作业问题。时至今日，这个问题只需要少于10行的代码便可以完成，普通的程序员都可以做到。

%---

2.3 深度学习基础---三种模型

2.3.1 起源

随着第一台电子计算机的问世，人工智能技术的发展拉开序幕。人工智能研究者基于对“智能“的不同理解，逐渐形成了符号主义、联结主义以及行为主义三大流派[17]。其中，人工神经网络[18]是连接主义的代表模型，它通过模拟生物神经网络及其机制，形成了由大量神经元通过互相连接作用形成的网络系统。从上世纪40年代的M-P神经元和Hebb学习准则，到50年代的感知器模型，再到60年代的自组织映射网络，许多神经计算模型都发展成信号处理、机器视觉、自然语言处理等领域的经典方法[19]。然而，由于当时该方法的不完备性，对深层神经网络的研究曾一度消弭[21][22]。2006年，Hinton指出具有隐藏层的网络具有优异的特征学习能力，提出了一种面向复杂通用学习任务的深层神经网络。自此，崭新的神经网络研究时代开启，深度学习的概念开始被提出。

 2.3.2 单层感知器模型

自从Cajal开创了神经元学说以来，许多神经元传递信号的方式在被逐渐发现。下图为一个简化的神经元模型：

图2.3.1 简化的神经元模型

神经元在无信号输入时保持抑制状态。而当树突和胞体接受的兴奋电位总和超过某个阈值时，它会进入兴奋状态，并经由轴突发出脉冲，刺激其他神经元。而神经元的M-P模型[6]正是经由这个原理而提出，如下图所示：

图2.3.2 M-P模型

	图中的$x_1,x_2,……,x_n$表示其他相连的神经元传递的输入信号，$w_ij$代表每个输入信号的权值，f称为激励函数，$y_j$表示当前神经元的输出，那么可以得到表达式：
%y_j=f(b_j+∑_(i=1)^n▒〖(x_i×w_ij))〗
其中$b_j$代表使该神经元激活所需要的阈值。这种通过逻辑角度来描述神经元的方法，为神经网络的理论研究开辟了道路。

2.3.3 多层感知器模型

单层感知机虽然在某些问题上有效，但是并不能解决诸如线性不可分问题等多种问题[23]。因此，多层感知机的概念产生：它的结构由输入层、隐藏层及输出层构成。下图是一个多层感知器的结构图：
	 
图2.3.3 多层感知机结构图

注意到，这里每两层之间的任意神经元都与另一层的所有神经元相连接。与单层感知器模型相同，每个输入信号都有一个对应的连接权值，而每一个神经元的输入为前一层所有神经元的输出的加权和。假设$x_m^l$代表着第l层第m个神经元的输入值，$y_m^l$和$b_m^l$分别为该神经元的输出值和偏置值，$w_im^(l-1)$代表该神经元与第i-1层第i个神经元的连接权值，那么有：
%x_m^l=b_m^l+∑_(i=1)^k▒w_im^(l-1) ×y_i^(l-1)
$y_m^l=f(x_m^l)$

虽然多层感知器模型解决了单层感知器的缺陷，但是却带来了新的问题，如下图所示，如果需要处理一个1000*1000像素的3通道RGB图片，即使在隐藏层只设置了1000个结点，那么待处理的参数也是多达1e9的量级。这仅仅只是一层，如果需要得到拟合度高的网络，那么待处理的节点个数将是一个天文数字。同时，这么多参数也可能会带来过拟合的问题。因此，亟需一个新的模型来克服这一缺陷。

2.3.4 卷积神经网络模型

2.3.4.1 感受野

1962年，科学家发现在视觉皮层中存在一系列构造复杂的细胞，它们只能识别到视觉输入空间的局部区域，因此被称为“感受野”[25]。下图为一个感受野模式图：

图2.3.4 感受野模式图
图中最上方的小圆圈就代表着众多感受野，光感受器将得到的信息整合并传递到神经元中，最后再经过进一步处理传到大脑。在识别图片时，这些细胞首先捕捉图片的边缘和轮廓，然后将其整合成一些完整的校组件，最后再将这些小部分最终整合成待识别的对象。下图为人脸识别上的一个例子：

图2.3.5 感受野

2.3.4.2 边缘提取[26]

	上文在提及感受野时，描述了这些细胞会首先捕捉图片的边缘。那么基于这种边缘提取的思想，也同样被运用到了神经网络中。以提取水平和垂直的信息为例，如下图所示：

图2.3.6 边缘提取

	由于图像可以用多个矩阵来表示，因此期望能用一种矩阵运算实现上图的提取效果，这便引入了“卷积”这一概念。首先，卷积操作定义为对图像矩阵和滤波矩阵做多次内积，如下图所示：

图2.3.7 卷积                       图2.3.8滤波矩阵

内积指的是两矩阵对应元素相乘，并将其求和放在结果矩阵的相应位置上。在提取运算过程中，图中的卷积核（即filter）从左到右，从上到下移动过整个图像矩阵，每移动一步分别计算原图中被卷积核覆盖的部分与卷积核的内积。右图中的这个6*6的图像矩阵有一个明显的垂直边缘，那么经过与图中所示的卷积核作卷积运算，在最终结果4*4的矩阵中，确实提取到了这个垂直边缘，读者可自行模拟这个运算过程。

2.3.4.3 卷积层[26]

以上介绍了感受野和边缘检测的概念，通过运用这些思想，卷积神经网络（Conventional Neural Network, CNN）就模仿了大脑识别图像的这一过程。

	首先，卷积神经网络的第一层被称为卷积层，与上述介绍的操作相对应。它的目的是为了提取输入的不同特征，前文中提到的垂直边缘只是众多需要被提取的特征之一。当然，正如大脑中视觉细胞的提取结构，某些卷积层只提取一些诸如边缘和线条等较低级的特征，更高层的网络能从低级特征中进一步提取更复杂的特征。在如此多层的提取中，提取到的特征矩阵会越来越小，这会带来边缘信息丢失等问题，举例来说：

图2.3.9 边缘信息丢失

如图为一个5*5的单通道图片，通常使用一个3*3大小的卷积核运算得到一个3*3大小的运算结果。如果重复这样的操作多次，那么图像的大小最后可能变成了1*1，而1*1的图像毫无特征可言，这并不是我们期望得到的结果。于是接下来引入了零填充的概念：
	零填充操作非常简单，在每次卷积结束后，在图像最外层加上若干层0值。注意到，由于0在内积运算中对结果不会造成影响，因此不会对图片增加额外的干扰信息。零填充操作如下图所示：

图 2.3.10 零填充
除此以外，卷积层中还有步长、多卷积核以及参数计算公式等内容，虽然其概念和计算并不复杂，但都过于细节，不在此作详细阐述。
2.3.4.4 池化层

	其次，经过卷积层提取后的特征图数量可能仍然很大，此时需要下采样，即池化操作。目的有二，其一如上文所述是为了减少特征图数，其二还可以提高特征图的鲁棒性，防止过拟合。通常在卷积层后增加一个池化层，一般通过取区域最值、平均值的方法进行采样。最常见的池化层是规模为2*2， 步幅为2，对输入的每个深度切片进行下采样。每个操作对四个数进行，下图所示为一个最大值池化：

图2.3.11 池化

值得注意的是，池化操作不需要进行学习，也不像卷积需要通过梯度下降进行更新。

2.3.4.5 全连接层

卷积层+池化层可以统一看成CNN的特征提取层，那么在提取完以后，最终将学习到的特征应用于模型需要的任务。这一层的操作主要有两步：第一步将所有的特征图进行扁平化（变为1*N的向量），第二步再接一个或多个全连接层，进行模型的学习。下图中的Fully connection部分即为全连接层：

图 2.3.12 全连接层

2.3.4.6 CNN小结

下图为卷积神经网络训练的基本结构。与传统机器学习方法相比，CNN的权值共享思想使得网络中需要训练的参数变少，降低复杂度的同时，减少了过拟合，更易于训练。到此，对神经网络基础的介绍结束，下一节将结合图像处理中的一个具体方向---目标检测，进一步介绍神经网络在图像处理方向上的发展史。

图2.3.13 卷积神经网络训练的基本结构

上图为卷积神经网络训练的基本结构。与传统机器学习方法相比，CNN的权值共享思想使得网络中需要训练的参数变少，降低复杂度的同时，减少了过拟合，更易于训练。到此，对神经网络基础的介绍结束，下一节将结合图像处理中的一个具体方向---目标检测，进一步介绍神经网络在图像处理方向上的发展史。

2.3.2 深度学习图像处理的发展---以目标检测为例[26][27][28][29]

2.4.1 起源

ImageNet Large Scale Visual Recognition Challenge（ILSVRC）是2010年到2017年间举办的机器视觉领域最受追捧且最具权威的学术竞赛之一。在图像分类这个竞赛分支中，2010年与2011年的获胜团队使用了传统的图像分类算法，最低取得了28.2 \%的错误率。而2012年的获胜团队首次将深度学习用于大规模图像分类，使得其错误率比亚军低了10 \%。从这一年开始，深度学习方法在图像领域开始被广泛采用，新的模型的涌现不断降低错误率，而大规模数据集的出现更使其得到了充分的训练。接下来本文会从目标检测方面更详细的介绍其发展。

2.4.2 R-CNN

物体检测相比于图像分类要来的更加复杂，因为我们需要对一张图像中可能含有的多个不同种类的物体进行定位和识别。2012年深度学习所取得的成功也同样影响到了物体检测的领域，Girshick等开始目标检测中运用卷积神经网络，提出了R-CNN模型。在介绍R-CNN之前，我们先简单介绍一下目标检测的滑动窗口方法。它的核心思想就是从左到右，从上到下地滑动窗口，利用分类识别目标。如下图所示：

图2.4.1 滑动窗口

但是这样的思想存在一个问题：如左图所示，识别人的窗口需要一个竖长方形窗口。但是如果把这个窗口平移到右边，试图用它来识别车辆，但是由于车辆的宽度大于长度，能符合这个窗口的概率不大。同样，如果用右图中识别车辆的横长方形识别人，也同样不成功。因此需要提前设计好多个可能会用到的窗口大小，每种窗口都滑动提取数张图片。提取成功后，将其输入CNN分类器中完成后续操作。可以发现，这是一种暴力式的操作，不仅会消耗大量的计算力量，由于人为设定的窗口大小大概率不准，因此其识别率也不高。这种思想的意义在于提供了一种解决目标检测问题的思路。

R-CNN模型框架如下图所示，它与上述方法不同，不使用暴力方法，而是用候选区域方法，改变了目标检测的模型思路。R-CNN模型在当时以优异的性能令世人瞩目。后续的SPPNet、Fast R-CNN、Faster R-CNN都是在此基础上加以改进的。

图2.4.2 R-CNN示意图

具体来说，第一步是生成候选框。这里用到了一种叫做选择性搜索（Selective Search）的方法。首先将每个像素作为单独的一组，然后计算每一组的纹理，将两个最接近的组结合起来，重复这样的操作直到所有区域不能再结合为止。下图展示了如何使区域增长：

图2.4.3 选择搜索算法

这一步中，我们在一张图片上提取了大约2000个候选区域。需要注意的是这些候选区域的大小不固定，而CNN需要固定大小的输入，所以需要再对候选区域做一些尺寸上的修改。在处理完毕后，第二步CNN将在这些候选区域的基础上，提取出更高级、更抽象的特征，作为下一步分类器的输入，提取的这些特征通常会储存在磁盘中。

图2.4.4 特征提取步骤

第三步，R-CNN使用支持向量机（SVM）进行二分类，这是整个流程中最重要的一部分。假如共需检测n个类别，那么需要n个不同类别的SVM分类器，并且每个分类器会对所有2000个候选区域的特征向量进行打分，判断其是否为当前这个SVM需要判断的物体类别。具体来说：

图2.4.5 SVM分类

如图所示，如果需要识别的类别为猫与狗，那么首先对这2000个候选区域是否为猫进行打分。注意到这里的分数只能为是或否，这就是所谓二分类器的意义。同理，狗类别的SVM对每张图片是否为狗进行打分。经过所有SVM判断结束后，我们可以大体了解这些候选区域中哪些是无关背景，哪些是需要的目标，以及其具体类别。

第四步，我们发现第一步中得出的建议框有重叠的情况，因此需要利用上一步的打分结果剔除重复的建议框。这其中又引出了一种新方法：非极大值抑制（NMS）。其首要目的就是一个物体只保留一个最优的框，去除那些冗余的候选框。首先，它根据上一步中的打分结果进行概率筛选，将2000个区域筛选至一个很小的数值。而这些剩下的框几乎都包含了一部分的图像内容，用上述方法得到的打分结果类似，难以再次筛选。因此下一步需要分别计算每一个候选框与实际目标区域的交并比（IOU），即两图像边框交集的面积与两图像边框并集的面积之比，以及各候选框两两之间的交并比。以下图为例：

图2.4.6 非极大值抑制

我们发现经过概率筛选后还剩下五个候选框，其中A、C框对应黑车，B、D、E对应白车。那么接下来就要从这两组中分别保留一个与对应车辆重合度最大的候选框。图中已经标识了它们各自相对于车辆的IoU，那么对于黑车而言，最高的是A框。再对A框与C框计算IoU，如果其IoU>0.5，那么意味着两者重合率较高，可以删除掉其中与目标物体重合率不高的，也就是C。对B、D、E而言这个步骤同理，最终只保留了B。最终在这五个候选框中检测出了A与B这两个目标。当然，在这个过程后依然还会有位置不准确的情况出现，因此需要用回归训练的方式修正这些候选区域。由于涉及到较多回归知识，在此不再过多赘述。通过以上四个步骤，R-CNN最终输出了图片中需要被检测的目标位置。

图2.4.7 VOC2010数据集运行结果
这种方法在RASCAL VOC 2010数据集中取得了出色的效果，mAP（对所有目标类别检测的效果取平均值）达到了53.7 \%。但是它仍存在一些尚未解决的问题：首先，由于在提取特征阶段，依然采用传统的选择搜索方法，导致候选区域的数量过多。而对于每个特征，所有的这些候选区域都将被输入到CNN中提取一次，不仅占用了大量磁盘空间，还非常耗时。在上述数据集中，检测一张图像的平均时间高达47s。其次，图像在统一尺寸的过程中易失真，导致不期望看到的几何形变。如下图所示，对这些问题的思考最终导致了SPP-Net的产生。

图2.4.8 SPP-Net引入 
2.4.3 SPP-Net[32]

	为了解决上述R-CNN的部分缺陷，He等在2014年提出了SPP-Net，主要创新点在于将一个空间金字塔池化层（spatial pyramid pooling，SPP）添加在卷积层和全连接层之间，从而能适应任何尺寸的图像输入。同时，它还对特征提取步骤进行了修改，使得不需要每个候选区域都经过CNN，只需要将整张图片输入，使运行速度得到了极大的提高，该层结构如下图所示：

图 2.3.9 SPP-Net

上一行的流程为标准R-CNN的流程，而下一行为将其改进后的SPP-Net流程。可以发现R-CNN先对图像进行选择与变换，再加入卷积层进行提取特征，而SPP-Net则直接将整张图片加入卷积层中，得到全图的特征图。随后，通过映射的方式将选择性搜索在原图中的位置对应到特征图中的相应位置。而这些映射过来的特征向量，经过空间金字塔池化层（spatial pyramid pooling，SPP），输出固定大小的特征向量，后续部分与R-CNN相同。那么下面将介绍每一步的细节。
对于映射操作，我们已经通过CNN提取到了整张图片的特征图，并通过选择性搜索算法得到了原图的候选框，接下来需要建立起两张图中对于同一物体的坐标对应关系。如下图所示：

图2.4.10 映射

这种映射已有明确的数学公式描述，在此不做深入探讨。

下一步，由于在上一步中映射到的特征图大小不一，但是对于一个固定的CNN，全连接层的输入是一个提前设定好的固定的数值，那么我们需要将特征图转换成固定大小特征向量。而这就是SPP的核心作用。假设需要被处理的某个候选区域的特征图的大小为12*10*256，那么如下图所示，SPP会通过池化的方式，把每一个候选区域都变成1*1、2*2和4*4的三张子图，连接到一起变为（16+4+1）*256=5376的大小，并将其交给下一步处理。可以发现，对每一个候选区域的特征图都可以转换为这种形式，这样就实现了大小的统一。

图 2.4.11 空间金字塔池化层
SPP-Net通过上述方式降低了卷积运算的数据量，减少了运算时间。然而，SPP-Net算法依然使用了R-CNN的框架，由于SVM的存在，特征需要存入磁盘，空间占用较大。

2.3.2.4 Fast R-CNN[33]
	经由R-CNN和SPP-Net的经验，在2015年，R-CNN的提出者之一Girshick又再次提出了Fast R-CNN模型。他发现，上述的两种方法都需要多次训练模型，如训练卷积神经网络和训练SVM。因此，他提出的RoI pooling整合了整个模型，将CNN、RoI pooling、分类器甚至前文中省略的回归模块一起训练。下面给出了Fast R-CNN的整体流程图：

图2.4.12 Fast R-CNN
首先，Fast R-CNN的前几个操作与SPP-Net相同。它首先将整个图片输入到一个基础卷积网络中，得到整张图的特征图，然后再将原图中选择性搜索算法的得到的候选框映射到特征图中。接下来，类似于SPP-Net中的SPP层，Fast R-CNN给出了一个RoI pooling层。可以发现它其实是一个简化版本的SPP层。如下图所示：

图2.4.13 ROI pooling层
在SPP层中，我们的得到了大小为1*1、2*2和4*4的三张子图，但是在RoI pooling层中我们仅仅通过池化的方式得到了一张4*4的子图。虽然前者相对而言较准确一些，但是后者的速度快很多。此外，Faster R-CNN的整个框架都直接用一个神经网络相连，这是因为R-CNN中CNN的训练与SVM的训练互相独立，两者的参数都无法互相影响，将SVM分类通过其他方法代替以后，使得所有过程都在整体的框架中进行，前后部分互相进行反馈与调整。并且，由于不再存在两个的独立的过程，中间也无需再用硬盘空间暂时保存信息，一切都在内存中存储，解决了磁盘占用大的问题。
在上述反馈过程中，需要进行损失的计算，这一部分同样涉及到了复杂的数学公式，同样不再赘述。
在同样的RASCAL VOC 2007数据集，Fast R-CNN的mAP达到了68.4\%。但是，无论是R-CNN，SPP-Net还是Fast R-CNN，其都是建立在传统的选择搜索的基础上，而这成为了上述方法的瓶颈。

2.3.2.5 Faster R-CNN[34]
为了彻底解决这个短板，同年，Ren等针对选择性搜索算法提取特征较慢的问题，提出了Faster R-CNN。在该论文中，区域建议网络（RPN）的概念首次被提出，它被设计用来替代传统的选择性搜索算法。简单来说，Faster R-CNN中也加入了一个提取边缘的神经网络，即寻找候选框的工作也交给神经网络了，如下图所示：

图2.4.14 R-CNN算法进化
那么需要关注的重点就在于其代替选择性搜索方法的区域建议网络，其主要作用也是为了得出比较准确的候选区域。它的思路是用n*n（默认3*3）大小的窗口去扫描特征图，并且在每个被扫描的位置中都考虑9种可能的参考窗口（anchors），如下图所示：

图2.4.15 区域建议网络
注意到，上述的9种参考网络非常考究。它使用了128、256、512三种尺度的窗口大小，对于每种尺度分别又有1：1、1：2、2：1三种长宽比，共计3*3=9种参考窗口。如下图所示：

图2.4.16 参考窗口形状与大小
在得到上述候选区域后，我们将它进行分类判别为前景或背景，并用回归方法修正位置，再加入RoI层中，后续的流程与Fast R-CNN类似。
	这种方法提出RPN网络取代了传统的选择搜索算法。自此，目标检测的四个步骤，即寻找候选区域、提取特征、目标分类、调整位置被统一到了一个深度神经网络框架中。

图2.4.17 Faster R-CNN结构
2.3.2.6 FPN[35]
然而，卷积神经网络在提取特征时，高层特征容易保留语义信息，低层特征容易保留位置信息。具体到目标检测领域，高层特征虽然能保留语义信息，但是由于此时特征图的尺寸较小，含有的位置信息并不多，不利于物体检测；低层网络虽然包含比较多的位置信息，但是图像的语义信息并不多，不利于图像的分类，这个问题在小尺寸物体检测上更为显著，这也就是为什么物体检测算法普遍对小物体检测效果不好的最重要原因之一。
基于这种问题，研究者们想到可以使用合并后的高层和低层特征来同时满足分类和检测的需求。传统的金字塔模型通过直接输入多尺度图像的方式构建多尺度的特征，如下图2.4.18所示，但是这样的计算成本太高。而R-CNN、Fast R-CNN和Faster R-CNN选择了只对卷积最顶层，即语义最丰富的一层进行特征提取，如下图2.4.19所示，但是如上文所言，这样对于目标的检测能力太差。其他方法如利用卷积网络的层次结构，在提取特征的过程中通过网络的不同层得到了多尺度的特征图。该方法虽然能提高精度，且基本上没有增加耗时，但是依然没有使用更加低层的特征图，对小物体的检测能力较弱。
    
图2.4.18	  单层特征						图2.4.19 多层特征
于是，综合以上几点考虑，Lin等提出了特征金字塔网络（FPN）方法。如下图2.4.20所示，FPN不仅使用了层次较深的特征图，并且浅层的特征图也被应用到FPN中。并通过自底向上，自顶向下以及横向连接将这些特征图高效的整合起来，在提升精度的同时并没有大幅增加检测时间。上述三个具体操作在此不作深入讨论，实验证明，该特征金字塔结构的特征提取效果优于单尺度的特征提取效果。

图2.4.20 FPN模型
到此，对双阶段目标检测算法的总结告一段落，列出下表再次将其进行对比：
表2.4.1 二阶段目标检测算法性能对比
算法	数据集	mAP(\%)
R-CNN	PASCAL VOC 2007	58.5
SPP-Net	PASCAL VOC 2007	59.2
Fast R-CNN	PASCAL VOC 2007	68.0
Faster R-CNN	PASCAL VOC 2007	78.0
FPN	MS-COCO	34.4

表2.4.2 二阶段目标检测算法综合对比
算法	创新点	优势	缺陷
R-CNN	将CNN应用于目标检测，利用选择分类算法和支持向量机	开创深度学习在目标检测领域的先河	耗时、占用空间大、易失真
SPP-Net	引入空间金字塔池化层，修改特征提取方式	解决失真问题，加快运行效率	占用空间依然很大
Fast R-CNN	引入softmax替代支持向量机，引入奇异值分解	减少了储存空间的占用	整体耗时受到选择搜索算法的限制
Faster R-CNN	引入区域建议网络，设计了多参考窗口机制	实现端到端的目标检测模型，降低用时	对小目标检测能力较差
FPN	引入多层特征与特征融合	大幅度提高对小目标的检测能力	

2.4.7 单阶段算法与两阶段算法[36]
上文中出现了两阶段这个字眼，两阶段指的是实现检测的方式主要有两个过程：第一步，提取候选区域；第二步，再对区域进行CNN分类识别。因此，它又被称为基于候选区域的检测。而单阶段目标检测算法没有候选区域提取阶段，它采用一个网络一步到位，训练过程也相对简单，可以在一个阶段直接确定目标类别并得到位置检测框。其代表算法有YOLO系列与SSD系列。由于篇幅问题与作者水平限制，本文中对单阶段目标检测算法不再作深入讨论。
2.4.8 总结

图2.4.21 目标检测路线图
上图给出了目标检测近20年来的里程碑算法，目标检测技术从传统的选择搜索算法到如今的深度学习算法，在提升精度的同时减少了对耗时和空间的需求。过去的几年中，不断有成熟的目标检测算法应用在工业界落地，不仅便利了人们的生活，也给学术界提供了丰富的灵感来源。下一小节，我们将实现一个简单的人脸识别技术用以演示。

2.4.9 深度学习图像处理举例
我们在github上下载了几个已经训练好的模型，并以人脸检测为例：
        
 “11.png” 和 “12.png”
使用如下的Python代码实现：
	检测的结果为：
%---
\subsection{1.3.成功案例}

长期以来机器学习总能完成其他方法难以完成的目标。例如，自20世纪90年代起，邮件的分拣就开始使用光学字符识别。实际上这正是知名的MNIST和USPS手写数字数据集的来源。机器学习也是电子支付系统的支柱，可以用于读取银行支票、进行授信评分以及防止金融欺诈。机器学习算法在网络上被用来提供搜索结果、个性化推荐和网页排序。虽然长期处于公众视野之外，但是机器学习已经渗透到了我们工作和生活的方方面面。直到近年来，在此前认为无法被解决的问题以及直接关系到消费者的问题上取得突破性进展后，机器学习才逐渐变成公众的焦点。这些进展基本归功于深度学习。

-苹果公司的Siri、亚马逊的Alexa和谷歌助手一类的智能助手能以可观的准确率回答口头提出的问题，甚至包括从简单的开关灯具（对残疾群体帮助很大）到提供语音对话帮助。智能助手的出现或许可以作为人工智能开始影响我们生活的标志。
-智能助手的关键是需要能够精确识别语音，而这类系统在某些应用上的精确度已经渐渐增长到可以与人类比肩[14]。
-物体识别也经历了漫长的发展过程。在2010年从图像中识别出物体的类别仍是一个相当有挑战性的任务。当年日本电气、伊利诺伊大学香槟分校和罗格斯大学团队在ImageNet基准测试上取得了28%的前五错误率[15]。到2017年，这个数字降低到了2.25%[16]。研究人员在鸟类识别和皮肤癌诊断上，也取得了同样惊世骇俗的成绩。
-博弈曾被认为是人类智能最后的堡垒。自使用时间差分强化学习玩双陆棋的TD-Gammon开始，算法和算力的发展催生了一系列在博弈上使用的新算法。与双陆棋不同，国际象棋有更复杂的状态空间和更多的可选动作。“深蓝”用大量的并行、专用硬件和博弈树的高效搜索打败了加里·卡斯帕罗夫[17]。围棋因其庞大的状态空间被认为是更难的游戏，AlphaGo在2016年用结合深度学习与蒙特卡罗树采样的方法达到了人类水准[18]。对德州扑克游戏而言，除了巨大的状态空间之外，更大的挑战是博弈的信息并不完全可见，例如看不到对手的牌。而“冷扑大师”用高效的策略体系超越了人类玩家的表现[19]。以上的例子都体现出了先进的算法是人工智能在博弈上的表现提升的重要原因。
-机器学习进步的另一个标志是自动驾驶汽车的发展。尽管距离完全的自主驾驶还有很长的路要走，但诸如[Tesla](http：//www.tesla.com/)、[NVIDIA](http：//www.nvidia.com/)、[MobilEye](http：//www.mobileye.com/)和[Waymo](http：//www.waymo.com/)这样的公司发布的具有部分自主驾驶功能的产品展示出了这个领域巨大的进步。完全自主驾驶的难点在于它需要将感知、思考和规则整合在同一个系统中。目前，深度学习主要被应用在计算机视觉的部分，剩余的部分还是需要工程师们的大量调试。

以上列出的仅仅是近年来深度学习所取得的成果的冰山一角。机器人学、物流管理、计算生物学、粒子物理学和天文学近年来的发展也有一部分要归功于深度学习。可以看到，深度学习已经逐渐演变成一个工程师和科学家皆可使用的普适工具。

\subsection{1.4.特点}

在描述深度学习的特点之前，我们先回顾并概括一下机器学习和深度学习的关系。机器学习研究如何使计算机系统利用经验改善性能。它是人工智能领域的分支，也是实现人工智能的一种手段。在机器学习的众多研究方向中，表征学习关注如何自动找出表示数据的合适方式，以便更好地将输入变换为正确的输出，而本书要重点探讨的深度学习是具有多级表示的表征学习方法。在每一级（从原始数据开始），深度学习通过简单的函数将该级的表示变换为更高级的表示。因此，深度学习模型也可以看作是由许多简单函数复合而成的函数。当这些复合的函数足够多时，深度学习模型就可以表达非常复杂的变换。

深度学习可以逐级表示越来越抽象的概念或模式。以图像为例，它的输入是一堆原始像素值。深度学习模型中，图像可以逐级表示为特定位置和角度的边缘、由边缘组合得出的花纹、由多种花纹进一步汇合得到的特定部位的模式等。最终，模型能够较容易根据更高级的表示完成给定的任务，如识别图像中的物体。值得一提的是，作为表征学习的一种，深度学习将自动找出每一级表示数据的合适方式。

因此，深度学习的一个外在特点是端到端的训练。也就是说，并不是将单独调试的部分拼凑起来组成一个系统，而是将整个系统组建好之后一起训练。比如说，计算机视觉科学家之前曾一度将特征抽取与机器学习模型的构建分开处理，像是Canny边缘探测[20]和SIFT特征提取[21]曾占据统治性地位达10年以上，但这也就是人类能找到的最好方法了。当深度学习进入这个领域后，这些特征提取方法就被性能更强的自动优化的逐级过滤器替代了。

相似地，在自然语言处理领域，词袋模型多年来都被认为是不二之选[22]。词袋模型是将一个句子映射到一个词频向量的模型，但这样的做法完全忽视了单词的排列顺序或者句中的标点符号。不幸的是，我们也没有能力来手工抽取更好的特征。但是自动化的算法反而可以从所有可能的特征中搜寻最好的那个，这也带来了极大的进步。例如，语义相关的词嵌入能够在向量空间中完成如下推理：“柏林-德国+中国=北京”。可以看出，这些都是端到端训练整个系统带来的效果。

除端到端的训练以外，我们也正在经历从含参数统计模型转向完全无参数的模型。当数据非常稀缺时，我们需要通过简化对现实的假设来得到实用的模型。当数据充足时，我们就可以用能更好地拟合现实的无参数模型来替代这些含参数模型。这也使我们可以得到更精确的模型，尽管需要牺牲一些可解释性。

相对其它经典的机器学习方法而言，深度学习的不同在于：对非最优解的包容、对非凸非线性优化的使用，以及勇于尝试没有被证明过的方法。这种在处理统计问题上的新经验主义吸引了大量人才的涌入，使得大量实际问题有了更好的解决方案。尽管大部分情况下需要为深度学习修改甚至重新发明已经存在数十年的工具，但是这绝对是一件非常有意义并令人兴奋的事。

最后，深度学习社区长期以来以在学术界和企业之间分享工具而自豪，并开源了许多优秀的软件库、统计模型和预训练网络。正是本着开放开源的精神，本书的内容和基于它的教学视频可以自由下载和随意分享。我们致力于为所有人降低学习深度学习的门槛，并希望大家从中获益。

\subsection{1.5.小结}

-机器学习研究如何使计算机系统利用经验改善性能。它是人工智能领域的分支，也是实现人工智能的一种手段。
-作为机器学习的一类，表征学习关注如何自动找出表示数据的合适方式。
-深度学习是具有多级表示的表征学习方法。它可以逐级表示越来越抽象的概念或模式。
-深度学习所基于的神经网络模型和用数据编程的核心思想实际上已经被研究了数百年。
-深度学习已经逐渐演变成一个工程师和科学家皆可使用的普适工具。

\subsection{1.6.练习}

-你现在正在编写的代码有没有可以被“学习”的部分，也就是说，是否有可以被机器学习改进的部分？
-你在生活中有没有这样的场景：虽有许多展示如何解决问题的样例，但缺少自动解决问题的算法？它们也许是深度学习的最好猎物。
-如果把人工智能的发展看作是新一次工业革命，那么深度学习和数据的关系是否像是蒸汽机与煤炭的关系呢？为什么？
-端到端的训练方法还可以用在哪里？物理学，工程学还是经济学？
-为什么应该让深度网络模仿人脑结构？为什么不该让深度网络模仿人脑结构？



2016年3月10日，AlphaGo登上了各大头条。谷歌研发的人工智能出乎大众意料，打败了围棋界公认的最优秀棋手之一李世石。
围棋与国际象棋一样，也是一种两人游戏。然而一直以来，围棋之于国际象棋，就像国际象棋之于井字棋。围棋更复杂，组合更多，也更难以预计。最优秀的围棋棋手有时候甚至会依靠某种直觉，但他们也无法解释这些直觉的来源---有些人甚至断言，这是一种超越机器能力的人类智能。长期以来，最优秀的算法也只能勉强追上中等水平的围棋棋手。

但在数年之间，某种机器学习的模型接连取得了不同凡响的成绩。物体识别、人脸识别、光学字符识别、自然语言处理、自动化翻译以及推荐系统从前都是机器无法解决的问题，但突然之间，它们全部都被深度学习解决了。硅谷的所有投资者挂在嘴边的就只剩下这几个词语，所有网络巨头都开始吹嘘自己现在能够提供些什么新服务。
深度学习那引发轰动的成功与经典软件工程的众多方法形成了鲜明的对比。人们习惯先构筑一个理论，用来保证某项技术能正确运作;然后，人们就会沮丧地发现技术超出了理论。这就引发了下面这个永恒的问题：理论和实践之间的差距从何而来?有一句俏皮话是这样说的：“理论上它们是一样的，但是在实践中嘛......”
但深度学习的情况似乎与此正好相反。似乎没有任何理论能够预言我们手头上的这项技术会取得成功。理论研究者被打了个措手不及。深度学习的性能简直好得离谱，但似乎没有人知道为什么。深度学习在实践中效果很好，但它在理论上是否可行?这次，实干家对这个问题感兴趣，是因为他们也感觉到，深度学习在理论理解上取得的任何进步，都有可能显著地甚至戏剧性地推动人工智能的发展。

特征学习
在深度学习大获成功之前，专业的数据科学家必须干预机器学习的算法，为的是对原始数据进行“清理”，预先简化数据的分析。深度学习研究的主要动机之一就是希望将数据预处理自动化，绕过对数据科学家技能的需求。这种做法让深度学习成为一把瑞士军刀，能够适应各种媒介，无论是图像、音频、视频、文字还是其他实时传感器的探测结果。

人工神经网络的灵感其实来自我们的大脑皮层。这是因为神经科学发现，我们的大脑会通过逐步抽象的方式来分析眼睛所看到的事物。眼睛中的传感器又被称为视锥细胞和视杆细胞，它们会探测那些令其进入激发状态的光线，得到光线的亮度和颜色。计算机科学家会说，这相当于图像中每一个像素的亮度和颜色。我们赋予深度神经网络的可观测变量的，正是这种原始数据。
然后，负责计算的第二层神经元一般会衡量相邻像素之间的相关度。人类的第二层神经元都会连接着眼睛的几个视锥细胞和视杆细胞。

人工智能的前沿就是所谓的卷积神经网络(convolutionneuralnetwork，简称CNN，因为它们与数学中的卷积运算有关)。这些网络的灵感就是来自视皮层的大体结构。图像分析的性能似乎依靠的正是这种抽象层次的堆积，而只有深度足够大的网络才能做到这一点。

单词的向量表示
人们在自然语言处理的经验中也发现了类似的现象。自然语言处理的难点之一就是含有大量的自然语言词汇。与其让单一的神经元专门负责每一个词语的理解，不如通过每个词语与其他词语之间的关系来理解它们。奇怪的是，用数学的语言来说，要做到这一点，我们可以将词语组成的空间嵌入高维线性空间中，这个维度一般来说在50和100之间。用不那么晦涩的语言来说，每一个词语都会对应着多个神经元的某个激活组合。
2013年，一个由托马斯·米科洛夫领导的谷歌研究团队成功让神经网络学会了之前说的这种英语单词的神经元表达，这是一项引人注目的成就。他们将这种表达称为word2vec。这种表达能以前述的方式将所有英语单词转换为高维空间中的向量。人们借此对所有单词组成的集合在表示方式上进行精简化与结构化，因为从信息的角度来看，单词的向量表示更紧凑，仅需更少的比特数就能描述。
但好处不止于此!单词的向量表示还能解释单词之间的众多联系。的确，米科洛夫研究团队的重大发现之一，就是单词之间的向量加法符合我们的直觉。比如，他们发现如果进行皇帝-男人+女人的向量运算的话，得到的结果大概就是皇后这个词语的向量表达[2]!
更令人着迷的是，我们在GAN中也能观察到这种现象。收集一些戴太阳镜的人的图像和没戴太阳镜的人的图像，你可以利用GAN计算出戴太阳镜的人对应的向量平均值以及没戴太阳镜的人对应的向量平均值。然后将这两个向量相减，得到的差就是某个向量，从直觉上来说，它对应着“戴太阳镜”。现在考虑一张没戴太阳镜的人的图像，计算它的向量表示，然后加上“戴太阳镜”的向量表示，就得到了一个新向量。现在利用GAN生成这个向量的对应图像。令人惊叹的是，得到的图像就是原来的图像中的人戴上了眼镜[3]!真是难以置信!
神经网络对不同抽象概念的向量求和结果的诠释竟然如此契合大脑对这些概念的直觉组合。



据此，能够阅读实体书的神经网络应该是一个深度神经网络，其中前几层是卷积神经网络，能够识别图像中的字符。然后，在负责视觉的这几层神经元之下，我们会发现一层能将字符合并为词语的神经元，以及另一层将词语转化为向量的神经元，就像word2vec所做的那样。最后，还有几层神经元负责解释那些表示词语的向量，将它们与其他概念联系起来。这些概念可以来自对图像或者其他信号的额外分析。
因此，高效的深度神经网络在看到猫的照片或者读到“猫”这个词语的时候，应该能够激活同一组神经元。为此，对原始数据的预处理似乎至关重要。一般来讲，对于原始数据的概括来说，无论是为了压缩体量宏大的原始数据，让计算能够在合理的时间内完成，还是为了揭示原始数据中相关性的切实解释，网络深度似乎都必不可少。


指数式的表达能力※
2016年6月16日，美国斯坦福大学、美国康奈尔大学以及谷歌大脑(GoogleBrain)的合作研究人员在arXiv上传了两篇非同凡响的论文。我被它们对抽象方法所获成功的别样解释深深迷住了。在本书中，我只能描述其中于我而言最有吸引力的那一篇。这篇论文题为《深度神经网络中经由暂态混沌出现的指数性表达能力》(“ExponentialExpressivityinDeepNeuralNetworksthroughTransientChaos”)。它结合了大量复杂、高深却鲜有联系的数学概念，比如混沌理论、平均场论和几何曲率。
我们慢慢解释。这篇论文一开头就认为，神经网络不过是一个复杂的数学函数。这个函数以一组测量数据作为输入，然后将其转化为深度概念的组合。测量数据的集合在数学上组成了一个所谓的向量，我们可以将它看成极高维度的空间中的一个点。而深度概念的组合也是如此，它本身就是极高维度的空间中的一个点。
这样的话，神经网络就可以被看成从第一个空间到第二个空间的几何变换。这篇论文提出的问题如下：“典型”的神经网络是如何将空间变形的?一个“随机选取”的神经网络平均而言是不是会将这些点到处移动?还是说，它会在变换中保留曲线的某些几何特性?
这篇论文的回答引人深思：几何结构的复杂度会随着神经网络深度的增加呈指数增长。更准确地说，论文研究了几何图形的全局曲率3。从直觉上来说，圆就是弯曲程度最小的闭合图形，因为它的弯曲是为了回到自己的出发点。此外，无论圆的大小如何，它的全局曲率总是等于$\tau\sim 6.28$。反之，一条在空间中的所有维度上都反复摆动的精细曲线会拥有非常大的全局曲率。（MW：傅里叶变换和曲率的关系：全局曲率越大，所需要的傅里叶级数越多，图形越复杂）
3全局曲率就是曲线方向上单位向量导数的范数的积分。
这篇论文证明了，对于足够“激动”的随机神经网络，几何图形的曲率会随着网络宽度的增长而呈多项式函数上升，并随着深度的增长呈指数上升。换句话说，比起宽度，网络的深度能让神经网络以更快的速度使对应的几何变换变得更复杂。
所以，这说明深度是识别与分析分形结构的关键，分形结构代表的就是那些并不总是光滑且正则的行为。然而，在我们周围的这个世界之中，分形结构似乎无处不在[4]，无论是在生物学、宇宙学还是金融之中!
同样，即使这种说法看似抽象甚至荒谬，但在所有图像组成的浩瀚集合之中，所有包含猫的图像组成的集合很可能同样有着分形结构。因此，网络的深度可能就是识别图像中有没有猫的关键，对于众多更复杂的任务来说也是如此。

\subsection{1.引言}

自古以来，人类对于自身大脑的探索从未停止。大脑，这个复杂而精密的器官，不仅赋予了我们思考、感知和学习的能力，更让我们在漫长的进化过程中脱颖而出，成为地球上最智慧的生物。然而，尽管科学家们已经取得了许多关于大脑结构和功能的重大发现，但大脑的工作机制仍然充满了未知和神秘。

在大脑的探索之旅中，一个尤为引人注目的领域是神经网络。神经网络，特别是人脑中的神经网络，是一个由数以亿计的神经元通过突触相互连接而成的庞大网络。这些神经元和突触不仅负责传递和处理信息，还通过复杂的交互作用，实现了我们的各种感知、思考和行为。

受到人脑神经网络的启发，科学家们开始尝试构建人工神经网络（Artificial Neural Networks，ANNs），以模拟大脑的信息处理过程。

人工神经网络 4大理论支柱为 “阈值逻辑”、“Hebb 学习率”、“梯度下降”、“反向传播”，前2个理论解决了单个神经元层面的建模问题，后2个理论则解决了多层神经网络训练问题。

2006年 Hinton 首次实现了5 层神经网络的训练，之后行业迎来爆发式发展，不断验证了该技术。

深度学习具备坚实的数学理论基础支撑。人脑绝大多数活动本质上都是广义计算问题，因此人脑其实是一个复杂的函数，深度学习就是去找到这个函数，万能近似定理则从数学上证明了一定条件下深度神经网络模型能够模拟任意的函数。
用于深度学习的算力以6年30 万倍的速度增加，算力是核心瓶颈也是未来提升的关键。从进化角度看，人的智能是一个随着神经元数量提升，从量变到质变的过程，这个量变过程对应人工智能中模型所用算力的提升过程。目前能够实现的 AI 模型中不论从神经元个数还是连接数量看，与真实人类还有较大的差距，未来随着现有芯片技术的不断推进，和突破冯诺依曼架构的类脑芯片等新技术的发展，算力的持续突破将会不断释放人工智能技术的潜力。


面对大数据的人类大脑

这些算法似乎都离日常生活非常遥远。你可能会说你能够一一记住自己接收的所有数据，但这就错得离谱了。

人们可能从未察觉，每时每刻充斥我们大脑皮层的数据量大得可怕，有如CERN需要处理的数据量。我们的视觉、听觉、嗅觉、触觉、温度感觉，以及其他数不清的感觉，每秒都会向我们传输大约十亿字节(1GB)的数据[4]。如果我们把数十年间收集到的数据累计起来的话，我们会发现感官接收到的数据总量需要用百亿亿字节(EB)作为计量单位!这可以说是大得可怕。

与现代信息技术面对的情况类似，我们的大脑也无法储存自身接收到的所有数据。它也不想这样做，因为其中绝大部分数据没有意义，它不得不忘记绝大部分数据。

在视觉数据的处理中，这种状况尤其突出。据神经科学家马库斯·赖希勒所言[5]，我们的视皮层拥有惊人的能力，可以将十亿字节的视觉数据转换为几千字节的有用数据，这种处理在每一秒都在进行，而且只需要消耗极少的能量[footnote视觉]。
[footnote视觉]:赖希勒估计，我们的视网膜每秒会收集到大约$10^10$比特的数据，但到达初级视皮层的第四层的数据只有大约$10^4$比特。


在更普遍的意义上，大脑会优先保留数据的“大体概念”，而不是具体细节。一般来说，我敢打赌你不能背出这本书前15章中的任意一句话。但在我的期望中，你应该记住了这些章节讨论的是贝叶斯公式、它的逻辑基础、它在归纳推理问题上的应用、它的历史、所罗门诺夫妖及其在博弈论中的应用，甚至还有它与演化理论之间的联系。即使你可能没有记住读到的任何一个句子，但我仍然希望你记住了读到的文字中各种基础概念的抽象表达。

对于接收到的信息只保留压缩过后的表示，这种能力并非弱点，而正是大脑的强大之处。我们一般都能够回忆起那些重要的事物，而且完全忘却那些对我们来说无足轻重的东西。

用贝叶斯帮助记忆
我们之前看到，人类大脑令人叹服的地方之一，就是它将极大量原始数据压缩为寥寥几个想法的能力。受此启发，人工智能研究者发明了
自编码器这一架构。
自编码器的任务，就是对大量信息进行精简，这些信息可以是高分辨率的图像或者整部电影，等等，最终只留下内容的精髓。换句话说，
这些神经网络尝试做到的，正是我们在语文课上被要求做的一种练习:撰写摘要。为了测试摘要的质量，人们也要求自编码器对自己生
成的摘要进行解压缩，设想与这一摘要相关的高分辨率图像或其他原始数据是什么样子的。
要做到这一点，关键之一就是利用贝叶斯方法。对数据重组，其实就是确定什么数据会得出眼前这个摘要。这正是贝叶斯公式的典型应
用!我们需要确定这个摘要的可能原因，因此

<!-- ![截屏2024-09-0418.50.53](/Users/Min369/Downloads/同步空间/Obsidian/书/截屏2024-09-0418.50.53.png) -->

另外，我们在下一章也会看到人工智能研究者如何利用这一公式，引入不常见的摘要，向机器赋予了创造性。

短期记忆与长期记忆

金特和朱莉娅·肖研究的是长期记忆。其中，信息被储存在大脑皮层这一神经网络的拓扑结构之中。问题在于访问这些数据可能很困难，需要测试信号在网络中传播的各种可能方式。回忆一首有名的歌曲的歌词一般来说就是这样的情况。在刚开始回忆时，我们通常想起几个词之后就想不起来了，就好像大脑中电信号的流动停滞了，或者流向了错误的方向一样。然而在重复这种流动之后，我们最终仍能找到正确的道路，回忆起整首歌的歌词。

人们在大型数据库中搜寻信息的时候也会碰到同样的问题。给谷歌和其他公司带来巨大财富的事物之一，就是为了加速信息搜索而整合互
联网数据的方法。但这并不只是信息整合的问题，人们也经常要在储存容量和存储速度之间对储存信息的媒介进行取舍。

递归神经网络

无论如何，人工神经网络的研究选择了第三条道路:神经信号的环路传播。包含这种环路的神经网络架构又叫递归神经网络。在处理当前数据时，这种网络能够利用前一时刻的部分信息。对于处理本质上有时序特性的数据，比如书中的文本或者讲话中的声音，递归神经网络已经成为最尖端的技术。

递归神经网络用到的技巧也包含在一类更广泛的拥有内部状态的算法之中。值得一提的是，这一类算法包括之前谈到的卡尔曼滤波器和隐马尔可夫模型。内部状态的任务就是以精练的方式提取出前一时刻的信息，用以更好地理解前一个被分析的数据。然而另一个问题就是，与数据读取相比，内部状态会以什么样的速度变化。
我有幸在自己学习数学和指导学生学习的过程中都观察到了一点东西。第一次读一份文献时(比如说你正在看的这本书)，最好不要停留在细节上，以免偏离阅读的主线。因为如果读得太慢，换种说法，就是对于短期记忆这一内部状态修改得太多的话，会使思绪偏离文献
想让我们思考的东西。因此，有时候先迅速浏览一遍，再多花点时间细细阅读，最后再速读一遍，这样可能效果更好。每一次阅读都会带来新的东西，因为每一次阅读都联系着短期记忆这一内部状态的不同变化动态。
到这里为止，我们讨论的还是对数据的线性阅读，但我们阅读或者聆听的方式并非完全线性的。的确，当某个句子非常晦涩的时候，我们一般会重读几次，然后才继续线性阅读。同样，有时候如果别人一句话还没说完，我们就无法理解。在梵语、印地语和日语中，动词位于句末，所以这种情况尤其显著。同样道理，在数学论文中，计算过程之后通常会有对计算的解释。还有些笑话、广告和电影，正是结尾决定了它们的意义。

为了同时利用过去和未来的数据理解现在，人们提出了另一种神经网络架构，那就是所谓的双向递归神经网络。为了将未来纳入其中，这些网络必须延长得出回应的时间。我们可以打赌，人类大脑中也有类似的结构。这也能解释我们在经过一段延迟之后才理解笑话的能力，甚至有时候这种延迟本身就引人发笑。

最后，人工智能领域的一项最新进展就是引入了迫使遗忘的神经元件。那就是所谓的LSTM架构，意思是“长短期记忆”(long-short-termmemory)。除了递归神经网络中常见的那些神经元环路之外，LSTM还拥有另一个额外的环路，它激活的时候会强制让所有神经元环路中的信号消失。这也许就解释了为什么我们在讨论中被打断之后，有时候很难回想起之前讨论的话题。

目前，在处理本质上有时序特性的数据方面，比如语音识别、自然语言分析，等等，LSTM及其变体似乎就是最尖端的人工智能方法。但我承认自己不理解它为何如此成功，我也承认自己没有足够的专业水平来预测这些领域未来的前沿进展。


\section{2.背景}



\subsection{第一阶段}

1943年，美国神经生理学家沃伦·麦卡洛克和数学家沃尔特·皮茨对生物神经元进行建模，首次提出了一种形式神经元模型。这个神经元模型通过电阻等元件构建的物理网络得以实现，被称为M-P模型。1958年，罗森布拉特又提出来了感知器，这意味着经过训练后，计算机能够确定神经元的连接权重。就这样，神经网络的研究迎来了第一次高潮。然而在1969年，明斯基等人指出感知器无法解决线性不可分问题，使得神经网络网络的研究陷入了第一次低谷。


\subsection{第二阶段}

1986年，Rumelhart(鲁梅尔哈特)等人提出了误差反向传播算法通过设置多层感知器，解决了线性不可分问题。从反向传播算法的发表开始，神经网络迎来第二次高潮。1989年，YannLecun等研究人员实现了一个7层的卷积神经网络LeNet-5，识别手写数字的精度达到了98%。在这段时间内，一大批学者和研究人员围绕Hopfield提出的人工神经网络方法展开了进一步研究，形成了20世纪80年代中期以来人工神经网络的研究热潮。可是好景不长，在20世纪末，支持向量机(SVM)和其他机器学习算法的流行程度逐渐高涨，更重要的是其他机器学习算法在实现效率上远超神经网络。随着PageRank等互联网新兴的算法不断被提出，学者们的研究重心也开始转移到其他机器学习算法上，神经网络的研究热潮又迎来了第二次低谷。



\subsection{3.第三阶段}

2006年，Hinton(辛顿)等人利用限制玻尔兹曼机对神经网络的连续层进行建模，使用逐层预训练的方法抽取模型数据中的高维特征，后来又提出了深度信念网络。这一技巧随后被众多学者推广到了许多不同的神经网络架构上，大大地提高了模型在测试集上的泛化效果。同时，随着硬件计算能力和算法的提升，网络隐层可以不断增加，于是以Hinton为代表的研究人员重新将人工神经网络定义为深度学习人工神经网络定义为深度学习，并进行推广。自201l年起，神经网络就在语音识别和图像识别基准测试中获得了压倒性优势，由此迎来了它的第三次高潮。





%![img](file:////Users/Min369/Library/Group%20Containers/UBF8T346G9.Office/TemporaryItems/msohtmlclip/clip_image002.jpg)









\subsection{3.结构}

![R-C](file:////Users/Min369/Library/Group%20Containers/UBF8T346G9.Office/TemporaryItems/msohtmlclip/clip_image004.jpg)

三个层次：输入层、隐藏层和输出层

输入层是神经网络的基层，它负责接收外界的输入数据。这些数据可以是图像、声音、文本等各种形式，经过预处理后转换成神经网络能够接受的形式。
隐藏层是神经网络的核心部分，它由多个神经元组成，负责处理输入数据并传递到下一层。隐藏层的设计和数量需要根据具体的任务来确定，它们的作用是对输入数据进行分类或预测。
输出层是神经网络的最高层，它负责将隐藏层处理后的结果转换成实际输出



圆形节点与人工神经元：

在人工神经网络中，每个圆形节点代表一个人工神经元。这些神经元通过特定的连接方式相互交互，模拟生物神经网络的工作原理。



连接与信号传递：

箭头表示从一个神经元的输出到另一个神经元的输入的连接。通过这些连接，信号可以在网络中传递，从一个人工神经元传递到另一个。



权重与激活函数：

每个节点都代表一种特定的输出函数，称为激活函数。每两个节点间的连接都有一个与之相关的权重值，表示前一个神经元对后一个神经元的影响程度。



网络输出：

网络的输出会根据网络的连接方式、权重值以及激励函数的不同而变化。通过调整这些参数，人工神经网络能够学习和适应不同的输入模式，产生预期的输出结果。


\subsection{4.如何训练}

4.1简介

反向传播算法（Backpropagation）是目前用来训练人工神经网络（ArtificialNeuralNetwork，ANN）的最常用且最有效的算法。其主要思想是：

（1）将训练集数据输入到ANN的输入层，经过隐藏层，最后达到输出层并输出结果，这是ANN的前向传播过程；

（2）由于ANN的输出结果与实际结果有误差，则计算估计值与实际值之间的误差，并将该误差从输出层向隐藏层反向传播，直至传播到输入层；

（3）在反向传播的过程中，根据误差调整各种参数的值；不断迭代上述过程，直至收敛。

4.2数学推导

1.变量定义

先定义一些变量：

%​![IMG_256](file:////Users/Min369/Library/Group%20Containers/UBF8T346G9.Office/TemporaryItems/msohtmlclip/clip_image005.jpg)表示第l-1层的第k个神经元连接到第l层的第j个神经元的权重；

%​![IMG_256](file:////Users/Min369/Library/Group%20Containers/UBF8T346G9.Office/TemporaryItems/msohtmlclip/clip_image006.jpg)表示第l层的第j个神经元的偏置；

%![IMG_256](file:////Users/Min369/Library/Group%20Containers/UBF8T346G9.Office/TemporaryItems/msohtmlclip/clip_image007.jpg)表示第l层的第j个神经元的输入，即：

%​![IMG_256](file:////Users/Min369/Library/Group%20Containers/UBF8T346G9.Office/TemporaryItems/msohtmlclip/clip_image008.jpg)

%![IMG_256](file:////Users/Min369/Library/Group%20Containers/UBF8T346G9.Office/TemporaryItems/msohtmlclip/clip_image009.jpg)表示第l层的第j个神经元的输出，即：

%​![IMG_256](file:////Users/Min369/Library/Group%20Containers/UBF8T346G9.Office/TemporaryItems/msohtmlclip/clip_image010.jpg)



%​![IMG_256](file:////Users/Min369/Library/Group%20Containers/UBF8T346G9.Office/TemporaryItems/msohtmlclip/clip_image011.jpg)其中表示激活函数。



2.代价函数

​代价函数被用来计算ANN输出值与实际值之间的误差。常用的代价函数是二次代价函数（Quadraticcostfunction）：

%​![IMG_256](file:////Users/Min369/Library/Group%20Containers/UBF8T346G9.Office/TemporaryItems/msohtmlclip/clip_image012.jpg)





​其中，x表示输入的样本，y表示实际的分类，
%![IMG_256](file:////Users/Min369/Library/Group%20Containers/UBF8T346G9.Office/TemporaryItems/msohtmlclip/clip_image013.jpg)表示预测的输出，L表示神经网络的最大层数。



3.公式及其推导

本节将介绍反向传播算法用到的4个公式，并进行推导。如果不想了解公式推导过程，请直接看第4节的算法步骤。

%![IMG_256](file:////Users/Min369/Library/Group%20Containers/UBF8T346G9.Office/TemporaryItems/msohtmlclip/clip_image014.jpg)首先，将第层第个神经元中产生的错误（即实际值与预测值之间的误差）定义为：





公式1（计算最后一层神经网络产生的错误）：



%​![IMG_256](file:////Users/Min369/Library/Group%20Containers/UBF8T346G9.Office/TemporaryItems/msohtmlclip/clip_image016.jpg)

其中，
%![IMG_256](file:////Users/Min369/Library/Group%20Containers/UBF8T346G9.Office/TemporaryItems/msohtmlclip/clip_image017.jpg)表示Hadamard乘积，用于矩阵或向量之间点对点的乘法运算。公式1的推导过程如下：



%![IMG_256](file:////Users/Min369/Library/Group%20Containers/UBF8T346G9.Office/TemporaryItems/msohtmlclip/clip_image019.jpg)

首先看上面证明的第一行，我们知道损失函数的形式，输出层的每一个输出都能够看作是损失函数的一个变量。

要是不理解的话，假设我们这里的损失函数为之前提到过的二次函数

在神经网络里面可以写为上面的形式，因为一个网络的输出就是最后一层的输出。这下能不能够明显的理解上面的那句话了呢？

那么证明第一行就是通过链式法则得到的。

然后只有%![IMG_256](file:////Users/Min369/Library/Group%20Containers/UBF8T346G9.Office/TemporaryItems/msohtmlclip/clip_image021.jpg)
和
%![IMG_256](file:////Users/Min369/Library/Group%20Containers/UBF8T346G9.Office/TemporaryItems/msohtmlclip/clip_image023.jpg)
是有依赖关系的，其他的没有依赖关系的求导为0不见了，因此，只剩下第二行。第二行出来了，又因为
%![IMG_256](file:////Users/Min369/Library/Group%20Containers/UBF8T346G9.Office/TemporaryItems/msohtmlclip/clip_image025.jpg)
，那么后面的也就都容易推出来了.





公式2（由后往前，计算每一层神经网络产生的错误）：

%​![IMG_256](file:////Users/Min369/Library/Group%20Containers/UBF8T346G9.Office/TemporaryItems/msohtmlclip/clip_image027.jpg)

推导过程：

%​![IMG_256](file:////Users/Min369/Library/Group%20Containers/UBF8T346G9.Office/TemporaryItems/msohtmlclip/clip_image029.jpg)



通过这个公式，只要我知道l+1层的误差，那么我就能够知道l层的误差，由此，我就能够倒退直到第1层的误差。
这个式子就体现了误差的反向传播的特点。
然后公式1和公式2结合起来，我们就能够求出神经网络中任何一层，任何一个神经元上面的误差。

公式3（计算权重的梯度）：

%​![IMG_256](file:////Users/Min369/Library/Group%20Containers/UBF8T346G9.Office/TemporaryItems/msohtmlclip/clip_image031.jpg)

推导过程：

%​![IMG_256](file:////Users/Min369/Library/Group%20Containers/UBF8T346G9.Office/TemporaryItems/msohtmlclip/clip_image033.jpg)

这个公式告诉我们，某个神经元上面误差对于偏置的偏导要等于在这个神经元上面的误差。而我们之前已经说过了，我们要求梯度的话就需要求出这些偏导，现在通过误差间接得到了偏导。很巧妙。

公式4（计算偏置的梯度）：

%​![IMG_256](file:////Users/Min369/Library/Group%20Containers/UBF8T346G9.Office/TemporaryItems/msohtmlclip/clip_image035.jpg)

推导过程：

%​![IMG_256](file:////Users/Min369/Library/Group%20Containers/UBF8T346G9.Office/TemporaryItems/msohtmlclip/clip_image037.jpg)

4.反向传播算法伪代码



输入训练集



对于训练集中的每个样本x，设置输入层（Inputlayer）对应的激活值：

前向传播：%![IMG_256](file:////Users/Min369/Library/Group%20Containers/UBF8T346G9.Office/TemporaryItems/msohtmlclip/clip_image039.jpg)

计算输出层产生的错误：
%![IMG_256](file:////Users/Min369/Library/Group%20Containers/UBF8T346G9.Office/TemporaryItems/msohtmlclip/clip_image041.jpg)

反向传播错误：
%![IMG_256](file:////Users/Min369/Library/Group%20Containers/UBF8T346G9.Office/TemporaryItems/msohtmlclip/clip_image043.jpg)

使用梯度下降（gradientdescent），训练参数：

%![IMG_256](file:////Users/Min369/Library/Group%20Containers/UBF8T346G9.Office/TemporaryItems/msohtmlclip/clip_image045.jpg)

%![IMG_256](file:////Users/Min369/Library/Group%20Containers/UBF8T346G9.Office/TemporaryItems/msohtmlclip/clip_image047.jpg)

4.3激活函数

激活函数（ActivationFunction），就是在[人工神经网络](https://baike.baidu.com/item/人工神经网络/382460)的神经元上运行的[函数](https://baike.baidu.com/item/函数/301912)，负责将神经元的输入映射到输出端，旨在帮助网络学习数据中的复杂模式。

常见类型：ReLU、Sigmoid、Tanh



Sigmoid激活函数的数学表达式为：

%​![IMG_256](file:////Users/Min369/Library/Group%20Containers/UBF8T346G9.Office/TemporaryItems/msohtmlclip/clip_image049.jpg)

导数表达式为：

%![IMG_257](file:////Users/Min369/Library/Group%20Containers/UBF8T346G9.Office/TemporaryItems/msohtmlclip/clip_image051.jpg)

函数图像如下：

%![IMG_258](file:////Users/Min369/Library/Group%20Containers/UBF8T346G9.Office/TemporaryItems/msohtmlclip/clip_image053.jpg)



tanh激活函数的数学表达式为：
%![IMG_256](file:////Users/Min369/Library/Group%20Containers/UBF8T346G9.Office/TemporaryItems/msohtmlclip/clip_image055.jpg)

函数图像如下：



%![IMG_256](file:////Users/Min369/Library/Group%20Containers/UBF8T346G9.Office/TemporaryItems/msohtmlclip/clip_image057.jpg)





实际上，Tanh函数是sigmoid的变形：

与sigmoid不同的是，tanh是“零为中心”的。因此在实际应用中，tanh会比sigmoid更好一些。但是在饱和神经元的情况下，tanh还是没有解决梯度消失问题。

什么情况下适合使用Tanh？

tanh的输出间隔为1，并且整个函数以0为中心，比sigmoid函数更好；

在tanh图中，负输入将被强映射为负，而零输入被映射为接近零。

Tanh有哪些缺点？

仍然存在梯度饱和的问题

依然进行的是指数运算





ReLU激活函数的数学表达式为：
%![IMG_256](file:////Users/Min369/Library/Group%20Containers/UBF8T346G9.Office/TemporaryItems/msohtmlclip/clip_image059.jpg)

函数图像如下：

%![IMG_257](file:////Users/Min369/Library/Group%20Containers/UBF8T346G9.Office/TemporaryItems/msohtmlclip/clip_image061.jpg)

什么情况下适合使用ReLU？

ReLU解决了梯度消失的问题，当输入值为正时，神经元不会饱和

由于ReLU线性、非饱和的性质，在SGD中能够快速收敛

计算复杂度低，不需要进行指数运算

ReLU有哪些缺点?

与Sigmoid一样，其输出不是以0为中心的

DeadReLU问题。当输入为负时，梯度为0。这个神经元及之后的神经元梯度永远为0，不再对任何数据有所响应，导致相应参数永远不会被更新

训练神经网络的时候，一旦学习率没有设置好，第一次更新权重的时候，输入是负值，那么这个含有ReLU的神经节点就会死亡，再也不会被激活。所以，要设置一个合适的较小的学习率，来降低这种情况的发生

4.1人工神经网络在各个领域中的应用



\subsection{人工神经网络在信息领域中的应用}

在处理许多问题中，信息来源既不完整，又包含假象，决策规则有时相互矛盾，有时无章可循，这给传统的信息处理方式带来了很大的困难，而神经网络却能很好的处理这些问题，并给出合理的识别与判断。

1.信息处理

现代信息处理要解决的问题是很复杂的，人工神经网络具有模仿或代替与人的思维有关的功能，可以实现自动诊断、问题求解，解决传统方法所不能或难以解决的问题。人工神经网络系统具有很高的容错性、鲁棒性及自组织性，即使连接线遭到很高程度的破坏，它仍能处在优化工作状态，这点在军事系统电子设备中得到广泛的应用。现有的智能信息系统有智能仪器、自动跟踪监测仪器系统、自动控制制导系统、自动故障诊断和报警系统等。

2.模式识别

模式识别是对表征事物或现象的各种形式的信息进行处理和分析，来对事物或现象进行描述、辨认、分类和解释的过程。该技术以贝叶斯概率论和申农的信息论为理论基础，对信息的处理过程更接近人类大脑的逻辑思维过程。现在有两种基本的模式识别方法，即统计模式识别方法和结构模式识别方法。人工神经网络是模式识别中的常用方法，近年来发展起来的人工神经网络模式的识别方法逐渐取代传统的模式识别方法。经过多年的研究和发展，模式识别已成为当前比较先进的技术，被广泛应用到文字识别、语音识别、指纹识别、遥感图像识别、人脸识别、手写体字符的识别、工业故障检测、精确制导等方面。

\subsection{人工神经网络在医学中的应用}

由于人体和疾病的复杂性、不可预测性，在生物信号与信息的表现形式上、变化规律（自身变化与医学干预后变化）上，对其进行检测与信号表达，获取的数据及信息的分析、决策等诸多方面都存在非常复杂的非线性联系，适合人工神经网络的应用。目前的研究几乎涉及从基础医学到临床医学的各个方面，主要应用在生物信号的检测与自动分析，医学专家系统等。

1.医学专家系统

传统的专家系统，是把专家的经验和知识以规则的形式存储在计算机中，建立知识库，用逻辑推理的方式进行医疗诊断。但是在实际应用中，随着数据库规模的增大，将导致知识“爆炸”，在知识获取途径中也存在“瓶颈”问题，致使工作效率很低。以非线性并行处理为基础的神经网络为专家系统的研究指明了新的发展方向，解决了专家系统的以上问题，并提高了知识的推理、自组织、自学习能力，从而神经网络在医学专家系统中得到广泛的应用和发展。在麻醉与危重医学等相关领域的研究中，涉及到多生理变量的分析与预测，在临床数据中存在着一些尚未发现或无确切证据的关系与现象，信号的处理，干扰信号的自动区分检测，各种临床状况的预测等，都可以应用到人工神经网络技术。

\subsection{人工神经网络在交通领域的应用}

今年来人们对神经网络在交通运输系统中的应用开始了深入的研究。交通运输问题是高度非线性的，可获得的数据通常是大量的、复杂的，用神经网络处理相关问题有它巨大的优越性。应用范围涉及到汽车驾驶员行为的模拟、参数估计、路面维护、车辆检测与分类、交通模式分析、货物运营管理、等领域并已经取得了很好的效果。

%![img](file:////Users/Min369/Library/Group%20Containers/UBF8T346G9.Office/TemporaryItems/msohtmlclip/clip_image063.jpg)

4.2人工神经网络应用和卷积对比

随着科技的快速发展，深度学习已经成为了人工智能领域的重要分支。在深度学习中，卷积神经网络（CNN）和人工神经网络（ANN）是两种最为常见的网络类型。它们在结构和应用上存在一些显著的区别。本节将对比分析CNN和ANN的区别，并通过实例阐述两者的应用。



CNN和ANN的区别主要体现在两个方面：结构和应用。首先，从结构上看，CNN侧重于处理图像数据，它的卷积层和池化层能够有效降低参数的数量，提高网络的性能；而ANN则没有这种结构上的限制，可以自由地设计神经网络结构来解决不同的问题。其次，从应用上看，CNN在处理图像任务时表现出色，如图像分类、目标检测等；而ANN则更加通用，可以应用于各种类型的数据和任务。
在应用案例方面，CNN和ANN各有优势。例如，在[图像识别](https://cloud.baidu.com/product/imagerecognition)任务中，CNN可以通过学习自动提取图像的特征，从而实现高效的图像分类。相比之下，ANN需要手动设计特征，因此在这方面的表现可能不如CNN。然而，在某些复杂的任务中，如[自然语言处理](https://cloud.baidu.com/product/wenxinworkshop)（NLP），ANN则可能更有优势。ANN可以通过构建深层神经网络，学习和理解语言的复杂模式和结构，从而实现高效的文本分类和语言翻译等任务。

%![img](file:////Users/Min369/Library/Group%20Containers/UBF8T346G9.Office/TemporaryItems/msohtmlclip/clip_image065.jpg)



猫是什么?
2012年，谷歌因一条奇怪而又令人不安的公告登上了各大新闻头条。
据这些头条所说，谷歌的人工智能发现了“猫”这个概念[9]!很多人可能觉得这是个平平无奇的新闻，但对我来说，这是机器学习的一个里程碑式的惊人突破，可能比AlphaGo在四年之后大败李世石给人留下的印象更深刻、更出人意料。
要理解这一点，首先要说明谷歌的人工智能是一个人工神经网络，带有一些接收器，让它可以“看见”数字化图像。谷歌给这个人工神经网络展示了1000万幅图像，其中不包含上下文。然后，谷歌将这个人工神经网络放进某种相当于核磁共振成像的处理系统中，用以测量它暴露在其他图像中的时候神经元的实时激活状态。谷歌意识到其中某些神经元大体上当且仅当向其展示的图像包含一只猫时才会激活!
真正给人留下深刻印象的是，这并不是谷歌在设计这个人工智能的时候设定的目标。这个人工智能的目标是以最合理有效的方式分析、处理并解释图像中的内容。人工智能必须为它看到的那些图像建立一个模型。为此，它必须做的事情之一就是解释图像中某些像素颜色之间反复出现的相关性。比如说，当图像中的某些像素构成了眼睛的形
状，那么在它的左边或者右边不远处通常会有一份与这些像素相似的复制品。一只眼睛的图像通常伴随着另一只眼睛的图像。怎么解释眼睛很少单独出现这一点呢?
最惊人的其实是，这个问题的解释似乎大体符合我们之中大多数人会给出的答案:因为照片拍到的动物通常有两只眼睛。人、狗和猫(几乎)都拥有两只眼睛这个事实对我们来说似乎如此微不足道，人们几乎不会对它花上任何时间进行思考。我们是如此习惯于日常生活中的事物，以至于我们甚至不再寻求解释---但这其实是个迷人的生物学问题，需要对视差的数学理解!
但对我来说更迷人的是，我们竟然可以(大体上)对于什么是人、什么是猫、什么又是狗达成共识。我们甚至没有意识到这些概念既抽象又模糊!猫是什么?怎么定义猫这个概念?更重要的是，猫存在吗?
要理解“猫”这个概念的奇怪之处，我们可以先强调这个定义有多么模糊。人们一般认为猫就是两只猫交配后产生的东西，换句话说，猫的双亲都是猫。到这里为止还没有什么惊人的东西，甚至这可能还是非常显然的东西。但其中大有深意。
比如说现在有一只猫，它的双亲都是猫，它的双亲的双亲也是猫，它的双亲的双亲的双亲也是猫，以此类推。但是这种对时间的回溯并没有尽头!根据这种推理，我们必须回溯生命演化树，到达猫还没有出现的年代!的确，如果我们回溯到几亿年甚至几十亿年前，我们会到达一个没有哺乳动物、脊椎动物，甚至连真核生物都不存在的时代!猫的双亲都是猫，或者说猫就是两只猫交配后产生的东西，这些说法在逻辑上就有矛盾[10]。
我预计到有些读者会提出用DNA的概念来定义猫这个概念。然而这会导致几个问题。首先，我们还没有对猫的所有基因组进行测序，而我们可能永远无法对所有猫的基因组进行测序，因为未来的猫现在还没有出生!所以怎么定义一组与猫对应的DNA代码并不是一个显然的问题。其次，即使我们成功做到了这一点，还要对动物进行基因组测序才能断定它是不是猫，这个解决方案在现实中完全不可行。再次，人们也许会严肃地考虑，包含猫的全套DNA的一个单独的猫细胞，比如猫毛中的细胞，它算不算一只猫?最后也是最重要的一点就是，DNA这个概念在大约半个世纪之前完全不存在。往好了说，这意味着薛定谔、达尔文和亚里士多德在谈论猫的时候其实不知道自己在说什么。
我们必须理解这个显而易见的事实:“猫”这个词并没有一个令人满意的严谨的定义。如果我问你猫是什么，你无法给出一个普遍适用、毫无争议的定义。而这不是什么惊人的事情，毕竟你不是通过这种方式学到怎么认知并使用“猫”这个概念的!如果你大概知道猫是什么，那是因为你已经看过上千张甚至上百万张猫的图像，你也听说过猫，也读过有关猫的材料。你通过观察大量数据学到了猫是什么。但你从来没有看见过关于“猫是什么”的形式定义!实际上，在人类历史中，没有人知道怎么给出“猫是什么”的形式定义。
还有更神奇的。必定存在第一个想到“猫”这个概念的人。因为这个人是第一个想到这个概念的，不会有人教他这一点，所以这个人就自己发明了“猫”这个概念。为什么?他又是怎么做到的?人类引入的新概念从何而来?这种能力是人类智慧特有的吗?
我觉得谷歌人工智能的发现非常震撼人心，因为它给出了这些问题的答案。不，这并非人类智慧特有的能力，因为谷歌的神经网络同样做到了这一点。而它做到这一点是因为它希望分析、综合并解释数字化图像中不同像素颜色之间的相关性，它希望得到一个模型来描述自己看到的东西，而它找到的模型自然导致了“猫”这个概念的诞生!


正确问题的近似解答的价值远远大于错误问题的精确解答的价
值。----约翰·图基(1915—2000)
真理(......)实在过于复杂，除了近似以外都不可行。--约翰·冯·诺伊曼(1903—1957)

无穷计算在数学里遍地都是。这些计算中最著名的可能是周角率(它与历史上“喧宾夺主”的圆周率[2]的关系是)。在14世纪，伟大的印度智者摩陀婆就发现了一个惊人的等式
的交错和的8倍!;也就是说，周角率不过就是奇数倒数
我们可能会认为这种无穷计算毫无用处。但事实上，无穷计算在应用数学的简单模型中无处不在，流体力学方程(导数和积分都是无穷步的计算结果)就是一个例子。这是因为，即使无穷计算只代表了一种不可计算的理想情况，但它一般能指出进行近似计算的巧妙方法。比如摩陀婆的等式就能用于计算的近似值，也就是说，它可以用于计算半径已知的圆的周长。工程师在实践中使用的正是这类近似方法。
如果你在常用的计算器或者谷歌上查询的值，它们很有可能会撒一点小谎，只给你提供小数点后十几位数字的近似值。这个问题并非只在上发生。你的计算器对于那些二进制表示并非有限的数都处理得非常糟糕[3]，这样的数包括无理常数，例如和，还有某些有理数，比如1/3和0.2。此外，因为计算器只能处理近似值，所以它可能会得出一些违反常理的结果，比如说，或者在比大得多的时候会得到，即使本身大于0。
数学家经常会认为，在计算机上进行的计算只不过是数学理论的近似。所罗门诺夫妖的立场却完全相反。对于所罗门诺夫妖来说，实数之类的对象只不过是一些让我们可以构建算法并对其更好地进行思考的模型。特别是对于尝试追随所罗门诺夫妖步伐的计算机科学家来说，他们用于尝试衡量置信度的模型并不是由实数作为参数组成的模型，而是那些储存在计算机文件中的模型，其中用到的只是数学模型中实数的截断。正因如此，与数学的理想状态相反，人工神经网络的大小可以用字节来计数4。

渐近展开
如果说对这样的数值进行高精度计算毕竟获益有限，那么对物理学、生物学或者数学中出现的曲线、函数和行为进行近似就会产生大量的应用。通过素数定理对素数计数函数进行近似就是这样的情况。
更一般的框架中存在这样一项工具，它对于任意模型都可以计算出比模型本身简单得多的近似版本。这项工具就是渐近展开。一般来说，渐近展开可以让我们用一条直线来逼近圆上的一小段圆弧，或者用平面来逼近(相当圆润的)地球的一小块表面。用代数术语来说，这相当于用形如的仿射方程来逼近所谓的非线性方程。对于那些变化足够小的现象来说，这种近似完全是可以接受的。这就解释了为什么我们要在学校里花上这么长的时间来研究这些简单的方程，以及为什么这些方程在科学中无处不在5。

对物理学家来说不可或缺的渐近展开，最精彩的例子可能就是爱因斯坦的广义相对论方程。它的渐近展开可以导出牛顿力学!换个说法，就是牛顿的力学定律，特别是关于万有引力的定律，只不过是爱因斯坦方程的近似，在有限的场景里完全适用。这种有限的场景就是所谓的“弱引力”情况6。[4]
6更准确地说，这应该是时空曲率极小的情况。

实用主义的限制
然而在现实中，计算能力是有限的。我们之前也看到了，在地球上执行超过$10^90$个计算步骤是种不切实际的愿望。这就大大降低了我们用贝叶斯公式进行计算的能力。

作为比较，人类大脑包含大约$10^15$突触，这意味着对大脑的完整描述至少需要比特的信息。研究所有包含这么多字符的理论，至少需要$2^{10^15}$次计算!因此，即使能用到古戈尔个宇宙，对这些理论应用贝叶斯公式也完全是一种幻想。

图灵的机器学习
1950年，艾伦·图灵将这个关于算法不可避免的复杂度的论证过程漂亮地转移到了人工智能问题上。在一篇发表在期刊《心灵:心理学与哲学季度评论》(Mind，aQuarterlyReviewofPsychologyandPhilosophy)上，题为《计算机制与智能》(“ComputingMachineryandIntelligence”)的论文中，图灵首先提出了一个与算法执行所需时间没什么关系的问题:“机器能思考吗?”图灵尝试回答的就是这个问题。然而“思考”这个词的模糊性让他尝试将问题精确化。与其考虑这个问题，图灵提出了另一个问题，就是机器是否能像人那样行动。

更准确地说，图灵提出了以下的测试:要求一位人类A与另外两个实体X和Y进行书面通信，而这两个实体X和Y分别是人类和一台机器。如果人类A无法分辨X和Y之中哪一个实体是人类的话，那么机器就通过了测试。换句话说，如果机器知道怎么模仿人类，而任何人类都无法分辨出这台机器不是人类的话，那么机器就成功通过了测试。这就是图灵所谓的“模仿游戏”，我们今天将它称为图灵测试[5]。

在论文的第3节到第5节中，图灵重提了他1936年的工作，严格定义了机器到底是什么---这个概念推动了计算机的发明，接着就是数字时代的到来!然后在论文的第6节中，图灵否定了尝试证明机器不可能思考的9个经典论证。但更重要的是，在论文的第7节，图灵预料到了其测试中的难点及解决方法。即使仍然不能说计算机当时已经存在，但图灵不仅已经预料到了它们会在未来出现，而且还预料到了它们以后的性能将足以在模仿游戏中取胜:“(工程上的进步)似乎不可能不足够(使其通过测试)。”对于图灵来说，“这个问题主要是编程问题”。
特别是，图灵以惊人的远见推测出，能成功通过图灵测试的程序代码至少需要大概$10^9$个字符来描述。也就是说，用我们在第7章引入的术语来说，图灵自己猜测，图灵测试的所罗门诺夫复杂度应该以十亿字节来计量。为了得到这个估计值，图灵依靠的是他所知的唯一一台能通过图灵测试的机器。对，我说的就是人类大脑!毕竟还有什么能比人类更善于模仿人类呢?对图灵来说值得庆幸的是，当时的神经科学已经给出了人类大脑复杂度的估计值，提出人类的神经元之间大约有$10^11$到$10^15$个突触---现代的估计值处于$10^14$和$510^14$之间。
图灵提出，只有一小部分突触对于通过图灵测试来说是必不可少的，这就是需要$10^15$个字符这个数字的来历。

同样出于他无比的智慧，图灵注意到人类大脑能够通过图灵测试。于是他猜测，通过模仿人类大脑成功完成图灵测试，我们能更好地构建能够通过测试的机器。然而图灵也注意到，儿童教育同样是大脑最终发育过程中的重要一环。“儿童的大脑可能类似于我们在文具店买到的笔记本，”图灵这样写道，“没什么机关，但有许多张白纸8。”图灵从而提出，为了帮助机器填满它自己的白纸，可以让它从数据中学习。学习机器的概念，也就是能够学习的机器，就此诞生。

因此，学习机器这个想法能够让机器自己写出包含数十亿字符的程序---必要时可以写得更长。用更贴近算法的语言来说，这相当于肯定了机器学习最终可以让我们研究并探索那些描述长度超过十亿比特的算法。而更关键的是，引导这种探索的并非程序员的指尖，而是原始数据，就像儿童接收到的那样。
需要特别指出，图灵的这段论证反驳了众多知识分子的想法，其中还包含一些专家，他们经常这样说:“机器学习效果不错---但数学家不知道为什么。”这是2015年《连线》(Wired)杂志上一篇文章的标题。然而早在1950年，艾伦·图灵就预言了机器学习未来会取得成功，甚至明确指出它会在20世纪末出现!更厉害的是，关于机器学习会在众多任务中超越人类编写的程序这一情况，图灵已经强调了它的确切原因:只有机器学习才能探索那些长度超过十亿字节的算法空间[6]。
同样，某些专家经常指责那些庞大的神经网络无法被明确解释。然而，巨大而高效的神经网络无法用寥寥几个字符描述出来，至少无法以准确的方式描述出来，这并不令人意外。这是因为，如果神经网络能被简单描述的话，这个不太长的描述就会是简短的一种算法，它能够生成另一个算法(神经网络)，并借此解决图灵测试之类的问题。
这样的话，这个简短的算法就能解决图灵测试。然而，我们刚才推测，这项测试无法被简短的算法解决。

现在剩下的就是明确指出用什么方法来探索长度巨大的算法组成的空间。然而艾伦·图灵并没有指出要做到这一点可以使用的方法，他只是断言学习能力将会是关键。在21世纪初，人们提出了大量的方法。我们之前已经谈到过其中几种。不分先后的话，我们可以举出线性回归、逻辑斯谛回归、决策树、决策森林、支持向量机、神经网络、贝叶斯网络，还有马尔可夫随机场。我们会在本章和后面几章中再详细讨论最后这三个机器学习算法。

实用贝叶斯主义
为了解决像图灵测试那样拥有巨大的所罗门诺夫复杂度的问题，实用贝叶斯主义者只能放弃贝叶斯公式。他更应该利用那些所需时间没有超出现实可能性的方法。由此我们需要调整“有用”这个概念。对于纯粹贝叶斯主义者来说，置信度高的理论就是有用的;但对于实用贝叶斯主义者来说，置信度很高，但所需计算时间超出常理的理论实际上并没有用。因此，实用贝叶斯主义者将注意力转向了所需计算时间较短的理论之中更可信的那一部分。

这就让我们能够解释素数定理给出的近似以及牛顿运动定律在实用中的置信度了。对于纯粹贝叶斯主义者来说，这两种描述都没有任何用处，因为两者分别被素数的精确计算以及爱因斯坦的广义相对论超越。然而实用贝叶斯主义者则乐于拥抱这些近似方法，因为它们对计算的需求明显小得多。计算的近似值需要的时间是的对数，而牛顿理论中的向量和微分计算与爱因斯坦理论中的张量计算相比远远更快。只要在我们身处的场景中，这两种近似精确到足以被接受(分别是值很大的情况和弱引力的情况)，那么对于实用贝叶斯主义者来说，这两种近似就比精确的版本更有用!
这种与计算时间有关的“有用”概念同样可以作为对神经网络惊人的成功的首要解释。这是因为，神经网络，或者至少是所谓的前馈神经网络，与那些更精细的算法相反，需要的计算时间不长，而且必定有所限制(图14.1)。实际上，如果我们考虑前馈神经网络的一个足够广泛的定义，那么这些网络组成的集合恰好就是快速并行算法的集合。因此，在前馈神经网络上进行机器学习，其实就相当于利用快速算法尽可能好地对数据进行解释(如果可能的话，也要利用合适的贝叶斯先验置信度)。所以，这就是迈向实用贝叶斯主义的第一步。
%<!-- /var/folders/gl/1x3lhvs57l72x96lt82xqx_r0000gp/T/TemporaryItems/NSIRD_screencaptureui_K6o907/截屏2024-09-0400.00.07.png -->

图14.1神经网络就是神经元之间一系列通信以及神经元本身进行的基本计算的总和。如果这种通信是无环的，我们就说这是个前馈神经网络，从直觉上来说，也就是通信只会朝一个方向进行，从来不会形成一个闭环ss